\documentclass[a4paper,12pt]{article}
\usepackage[utf8]{inputenc}
\usepackage[russian]{babel}
\usepackage[OT1]{fontenc}
\usepackage{amsmath, amssymb}
\usepackage{geometry}
\geometry{left=2cm,right=2cm,top=2cm,bottom=2cm}

\begin{document}

\tableofcontents

\newpage
\section{Билет №1. Феномен знания. Традиционное определение и его стандартный анализ}

\subsection{Ответ}

Аналитическая эпистемология рассматривает феномен знания прежде всего через процедуру концептуального анализа — метода, направленного на выявление необходимых и достаточных условий применения понятия путем исследования его логической структуры и интуитивных случаев употребления. Как отмечает Гаспаров И. Г., в англоязычной традиции эпистемология тождественна теории знания ($theory$ $of$ $knowledge$). Отправной точкой является поиск необходимых и достаточных условий, при которых субъекту $S$ можно атрибутировать знание пропозиции $p$.

Традиционное определение знания, восходящее к диалогам Платона «Менон» и «Теэтет», квалифицирует знание как «обоснованное истинное убеждение» ($justified$ $true$ $belief$, JTB). В «Теэтете» Платон последовательно рассматривает и отвергает определение знания как простого восприятия и как истинного мнения, приходя к формулировке знания как истинного мнения, сопровождаемого логосом (обоснованием). Хотя сам Платон в итоге выражает сомнения в полноте этого определения, именно эта трехкомпонентная структура стала основой для последующей эпистемологической традиции. Согласно анализу, представленному Лебедевым М. В. и Черняком А. З., это определение распадается на три обязательных компонента.

\begin{enumerate}
\item \textbf{Условие убеждения (ментальное условие).} Субъект $S$ должен принимать $p$, то есть находиться в определенном когнитивном состоянии относительно данной пропозиции. Гаспаров уточняет, что в аналитическом контексте термин $belief$ переводится как «убеждение» и трактуется максимально нейтрально: как то, что человек считает истинным, независимо от эмоциональной окраски или степени уверенности. Знание без убеждения невозможно, так как нельзя «знать» то, во что ты сам не веришь.

\item \textbf{Условие истины (алетическое условие).} Пропозиция $p$ должна быть фактически истинной. Эпистемология разграничивает знание и мнение именно по этому критерию: можно иметь ложное мнение, но нельзя иметь «ложное знание». Если субъект убежден в том, что Земля плоская, и имеет для этого аргументы, мы не скажем, что он «знает» это, так как само утверждение ложно. Знание каузально и логически связано с фактическим положением дел в реальности.

\item \textbf{Условие обоснования (эпистемическое условие).} Это условие призвано исключить «эпистемическое везение». Гаспаров указывает, что истинное убеждение может быть результатом случайного угадывания. Например, если человек ставит на номер в рулетке, веря, что он выиграет, и номер действительно выпадает, его убеждение истинно, но оно не является знанием, так как не имеет под собой рациональных оснований в момент формирования. Обоснование ($justification$) придает истинному убеждению статус знания, связывая когнитивные способности субъекта с истинностью пропозиции через логические доводы, свидетельства или надежные когнитивные процессы.
\end{enumerate}

Традиционный анализ утверждает, что знание имеет место тогда и только тогда, когда выполняются все три условия одновременно. Однако данная стандартная модель была дестабилизирована в середине XX века, что привело к современным поискам четвертого условия знания.

\subsection{Обзор источников}

Базовым источником для данной темы выступает работа \textbf{Гаспарова И. Г. «Knowledge»}. Глава 9 «Аналитическая эпистемология» в пособии Лебедева и Черняка дает системный взгляд на предмет. Структура работы Гаспарова (раздел 9.1 и 9.2) построена как историко-аналитическое введение: от концепции Бертрана Рассела о различии знания-знакомства и знания истин к классической JTB-модели. Гаспаров последовательно раскрывает каждый компонент JTB, объясняя, почему истины и убеждения недостаточно по отдельности. \textbf{Лебедев М. В. и Черняк А. З.} дополняют это строго семантическим описанием пропозициональных установок. \textbf{Анкин Д. В.} в своей «Теории познания» (раздел 4) фокусируется на логической структуре определения, объясняя разницу между необходимыми и достаточными условиями через формальные символы ($S$ знает, что $p$).

\subsection{Контрольные точки с подробными ответами}

\begin{enumerate}
\item \textbf{В чем заключается суть концептуального анализа в эпистемологии?}\\
Концептуальный анализ направлен на выявление необходимых и достаточных условий применения понятия «знание» через исследование его логической структуры и интуитивных случаев употребления.

\item \textbf{Какова роль диалогов Платона в формировании JTB-модели?}\\
В диалоге «Теэтет» Платон последовательно рассматривает определение знания, приходя к формулировке знания как истинного мнения с логосом (обоснованием), что заложило основу трехкомпонентной структуры JTB.

\item \textbf{В чем заключается суть «эпистемического везения»?}\\
Это ситуация, когда субъект обладает истинным убеждением, но связь между его верой и истиной случайна. Такое убеждение не считается знанием, так как оно получено не через рациональный метод.

\item \textbf{Можно ли знать то, во что не веришь?}\\
Нет, согласно стандартному анализу, убеждение является необходимым условием знания. Знание — это подмножество убеждений субъекта.

\item \textbf{Является ли истинное убеждение достаточным для знания?}\\
Нет, оно является только необходимым. Для статуса знания требуется третий компонент — обоснование, исключающее случайность.

\item \textbf{Какова роль JTB-модели в истории философии?}\\
Она служила доминирующим стандартом определения знания на протяжении почти двух с половиной тысяч лет, утверждая рационалистический идеал познания как процесса, подконтрольного разуму.
\end{enumerate}

\newpage
\section{Билет №2. Классическое определение истины, основные интерпретации и проблемы}

\subsection{Ответ}
Классическое определение истины базируется на концепции соответствия или корреспонденции. Согласно этому подходу, признаваемому еще Аристотелем, истина представляет собой отношение между суждением (утверждением о мире) и внешней объективной реальностью. Механика этого определения подразумевает, что высказывание является истинным, если то положение дел, которое оно описывает, действительно существует в мире. Например, фраза о том, что снег бел, оказывается истинной в том и только в том случае, если физический объект снег обладает свойством белизны. Здесь первичным носителем истинностного значения выступает предложение или пропозиция, а фактором истины становится конкретный факт.


Под «аналитической традицией» уместно понимать линию философии, где философские проблемы проясняются через \textbf{логический и семантический анализ языка} (лингвистический поворот, связь философии с логикой и семантикой).

\begin{enumerate}
  \item \textbf{Структурная (логико-атомистическая) интерпретация соответствия.}\\
  Рассел видит соответствие  прежде всего на уровне \textbf{базисных} (простейших) предложений: они истинны благодаря отношению к некоторому \textbf{событию/факту}, а истинность остальных предложений зависит от их \textbf{синтаксических отношений} с базисными. {Сложные предложения} получают истинность \emph{производно} — через логико-синтаксические связи с базисными. Поэтому истинность зависит от разложения: чтобы оценить сложное, мы должны знать истинность его базисных компонентов и структуру связок.
  Здесь речь не о физическом сходстве мысли и вещи, а о зависимости истинности сложного от логической формы и истинности базисного.

  \item \textbf{Семантическая (референциальная) интерпретация соответствия.}\\
  Здесь акцент на \textbf{семантических правилах языка}, то есть на отношениях между выражениями и тем, к чему они относятся. Тарский прямо определяет семантику как дисциплину об отношениях между \emph{выражениями} и \emph{объектами/положениями дел}, вводя типичные семантические понятия: \textbf{обозначение (designation)} и \textbf{выполнимость (satisfaction)} (Тарский, \S 5). На этой базе он строит определение истины через выполнимость: предложение истинно, если оно выполняется (выполнимо) всеми объектами/последовательностями в соответствующей модели (Тарский, \S 11).\\
  В более простом пересказе (с сохранением смысла): 
  \begin{enumerate}
    \item имена (единичные термины) получают \textbf{референцию} — «указывают» на объекты;
    \item предикаты получают условия \textbf{выполнимости} — какие объекты «подходят» под предикат;
    \item истинность предложения определяется тем, выполняются ли его условия в мире/модели.
  \end{enumerate}

\end{enumerate}

Проблемы классического подхода:
\begin{enumerate}
  \item \textbf{Неясность «соответствия».} Если понимать соответствие как «сходство», возникает псевдопроблема: как нематериальное высказывание «похоже» на материальный объект? Корректнее трактовать соответствие как семантическое отношение и/или как зависимость истинности от базисных фактов и синтаксиса.
  \item \textbf{Круг с «фактами».} В логической форме корреспондентной теории факт часто вводится как то, что делает предложение истинным; тогда возникает опасность объяснять истину через «то, что делает истинным истину». Эту трудность Рассел фактически фиксирует, разводя две формы корреспондентной теории (эпистемологическую: через опыт; и логическую: через факты) и показывая их различие по статусу закона исключённого третьего (Рассел, гл. XXI).
  \item \textbf{Парадоксы семантически замкнутого языка (антиномия лжеца).} Тарский показывает, что если язык способен говорить о собственной истинности (семантически замкнут), то при сохранении обычной логики возникают противоречия вроде «это высказывание ложно»; поэтому требуется различение \textbf{объектного языка} и \textbf{метаязыка} и иерархия уровней.
\end{enumerate}

\subsection{Обзор источников (привязка к ответу)}
\begin{enumerate}
  \item \textbf{А. Тарский, «Семантическая концепция истины и основания семантики».}\\
  Использовано для: (a) связи с аристотелевской интуицией истины как соответствия (Тарский, \S 3); (b) Т-схемы и требования материальной адекватности (Тарский, \S 4); (c) семантики как отношений «выражение — объект/положение дел», понятий обозначения и выполнимости (Тарский, \S 5); (d) запрета семантически замкнутого языка и различения объектного языка/метаязыка как решения парадокса лжеца (Тарский, \S 7–9); (e) идеи определения истины через выполнимость (Тарский, \S 11).
  \item \textbf{Б. Рассел, «Исследование значения и истины», гл. XXI «Истина и верификация».}\\
  Использовано для: (a) четырёх типов теорий истины и выбора корреспондентной (Рассел, гл. XXI); (b) различения «эпистемологической» и «логической» форм корреспондентной теории (Рассел, гл. XXI); (c) роли базисных (атомарных) суждений, соответствия событиям/фактам и производной истинности через синтаксические отношения (Рассел, гл. XXI); (d) обсуждения закона исключённого третьего в зависимости от выбора формы корреспондентной теории (Рассел, гл. XXI).
  \item \textbf{Я. Хинтикка, «Проблема истины в современной философии» (1996).}\\
  Использовано для: (a) рамки «язык как универсальный посредник» vs «язык как исчисление» и связи с выразимостью семантики (Хинтикка, раздел о контрасте двух взглядов); (b) тезиса о том, что тарсковский негативный результат зависит от предпосылок о логике, и обсуждения альтернатив (IF-логика) как способа переосмысления «невыразимости истины» (Хинтикка, фрагменты о Тарском и IF-логике).
  \item \textbf{Д. В. Анкин, «Теория познания» (2019).}\\
  Использовано для: (a) характеристики аналитической традиции через лингвистический поворот и логическую семантику (Анкин, гл. 2); (b) систематизации тем «корреспондентная концепция и её проблемы», «разграничение определения и критерия», «истина и язык» как структуры билета (Анкин, гл. 3).
\end{enumerate}

\subsection{Контрольные точки (кратко)}
\begin{enumerate}
  \item \textbf{Что такое «семантические правила» в ответе?}\\
  Это правила обозначения и выполнимости: как выражения соотносятся с объектами/условиями, при которых предложение истинно (Тарский).
  \item \textbf{Почему проблема лжеца опасна?}\\
  Потому что семантически замкнутый язык, в котором можно говорить об истинности собственных предложений, при обычной логике порождает противоречия; нужно различать объектный язык и метаязык (Тарский).
\end{enumerate}

\newpage
\section{Билет №3. А. Айер о специфике формальных наук}

\subsection{Ответ}

Айер в своей работе рассматривает статус логики и математики с позиций радикального эмпиризма. Основная трудность для эмпирика заключается в объяснении необходимых истин. Если верить в то, что всё знание черпается из опыта, то возникает противоречие, так как опыт никогда не дает логической гарантии универсальности. Юм показал, что факт подтверждения закона в $n-1$ случаях не гарантирует его выполнение в $n$-м случае. Следовательно, утверждения формальных наук должны быть либо эмпирическими обобщениями, либо чем-то иным, не относящимся к описанию фактов.

Эмпирист Милль пытался доказать, что истины логики — это индуктивные обобщения с очень высокой вероятностью. Айер этот взгляд отвергает. Он указывает, что математические пропозиции не могут быть опровергнуты опытом. Если при пересчете пяти пар объектов мы обнаруживаем девять предметов, мы не ревизуем математическую истину $2 \times 5 = 10$. Вместо этого мы выдвигаем эмпирическую гипотезу о том, что один объект слился с другим или мы ошиблись в счете. Это доказывает, что истины формальных наук обладают иным статусом: они необходимы, потому что мы не позволяем им быть ложными.

Специфика формальных наук по Айеру заключается в их аналитическом характере. Это тавтологии. Аналитическая пропозиция истинна исключительно в силу определений составляющих её символов. В отличие от синтетических суждений, которые сообщают информацию о внешнем мире и верифицируются через опыт, аналитические суждения лишены фактуального содержания. Они не говорят ничего о реальности, но фиксируют наши лингвистические соглашения. Например, утверждение о том, что сумма углов евклидова треугольника равна 180 градусам, является следствием определений самой геометрии Евклида.

Айер утверждает, что формальные науки полезны, несмотря на их тавтологичность, из-за ограниченности человеческого разума. Если бы наш интеллект был бесконечно могуч, мы бы видели все следствия определений мгновенно. Дедуктивный вывод не дает приращения знания о мире, но позволяет сделать явным то, что уже было имплицитно заложено в посылках. Таким образом, логика и математика — это правила тавтологического преобразования выражений, которые служат проводниками в эмпирическом поиске, не добавляя к нему новой фактуальной информации.

\subsection{Обзор источников}

Глава о знании a priori в книге Айера Язык, истина и логика построена как последовательное обоснование эмпиристской трактовки математики. Автор начинает с критики рационализма, который считает мышление надежным источником внеопытного знания. Затем Айер анализирует теорию Милля и доказывает её несостоятельность через примеры с пересчетом объектов. Основная часть главы посвящена пересмотру кантовского разделения суждений. Айер уточняет критерий аналитичности, избавляя его от психологизма Канта и заменяя его логико-лингвистическим критерием. Завершается работа объяснением того, почему тавтологии могут быть неожиданными и как они функционируют в дедуктивных системах вроде тех, что описаны в Principia Mathematica.

\subsection{Контрольные точки с подробными ответами}

1. Почему Айер считает истины логики лишенными фактуального содержания?
Потому что они не описывают свойства объектов мира, а лишь фиксируют правила использования языка. Отрицание тавтологии ведет к самопротиворечию, а не к противоречию с фактами.

2. В чем заключается ошибка Милля в трактовке математики согласно Айеру?
Милль считал математику индуктивным обобщением. Айер же показывает, что математические истины неопровержимы опытом: любой негативный опыт мы интерпретируем как ошибку измерения или наблюдения, а не как ложность математики.

3. Какова роль ограниченности разума в существовании формальных наук?
Если бы разум видел все следствия из определений сразу, логика и математика были бы не нужны. Они полезны только потому, что помогают нам вычислять сложные следствия наших собственных лингвистических соглашений.

4. Как Айер определяет аналитическое суждение?
Это суждение, истинность которого зависит исключительно от определений входящих в него символов. Оно остается неизменным вне зависимости от того, какова реальность.


\newpage
\section{Билет №4. К. Поппер о специфике формальных наук}

\subsection{Ответ}

Поппер рассматривает специфику формальных наук через проблему их \textit{применимости} к реальности: почему логика и математика, будучи абстрактными символическими построениями, оказываются полезными в описании материального мира. Ключевая мысль здесь в том, что формальные системы сами по себе не являются описаниями фактов.

Пока мы имеем дело с логикой или арифметикой как с формальной системой (с символами, формулами и доказательствами), получаемые утверждения обладают характером логической истинности: они выводимы в силу принятых правил и определений и потому не зависят от того, что происходит в природе. Однако ситуация меняется, когда мы начинаем \textit{применять} эту формальную систему к миру, то есть читать её выражения как утверждения о физических предметах, измерениях и событиях. Тогда появляется фактическое содержание, а вместе с ним — и риск ошибки.

Отсюда центральный тезис Поппера: \textit{в той мере, в какой математическое утверждение применимо к реальности, оно перестаёт быть безусловно истинным «в силу формы» и становится опровержимым}. И наоборот: если мы настаиваем, что оно истинно при любых обстоятельствах и не допускает опровержения опытом, то тем самым мы удерживаем его в сфере абстрактного формального построения, а не описания реальности.

Граница хорошо видна на примере арифметики. Внутри самой арифметики равенство $2+2=4$ выступает как следствие определений и правил обращения со знаками, то есть как формальная (логическая) истина. Но при применении к физическим объектам мы неизбежно добавляем содержательные предпосылки: что именно считаем «двумя объектами», что признаём сохранение объектов, по какому критерию считаем их теми же самыми и т.\,п. Поэтому в реальных ситуациях наивное «сложение» может не сработать: физические процессы (вроде исчезновения, появления, деления, взаимодействия живых существ) нарушают те допущения, на которых держится простой пересчёт. В таком случае под угрозой оказывается не формальная истинность $2+2=4$ как элемента арифметики, а корректность \textit{конкретного применения} арифметики к данной эмпирической ситуации.

Этот вывод Поппер подкрепляет своим пониманием роли логики. Следуя Тарскому, он принимает семантическое понимание истины как соответствия фактам, но подчёркивает: логические истины сами по себе не сообщают фактов о мире. Их функция иная: они задают условия корректного следования, то есть гарантируют сохранение истинности при переходе от посылок к заключению. Поэтому логика и математика выступают в науке не как самостоятельные описания реальности, а как инструменты контроля рассуждений и критической проверки: фактическое содержание вносится не ими, а нашими интерпретациями и эмпирическими предпосылками, которые и могут быть опровергнуты опытом.

Итог: формальные науки остаются «чистыми» и нефактуальными, пока рассматриваются сами по себе; но при применении к реальности они начинают работать как часть эмпирического знания — и именно тогда попадают в область проверяемости и возможного опровержения.

\subsection{Обзор источников}

Опора на обсуждение применимости логики и арифметики к реальности у Поппера (в «Предположениях и опровержениях», в полемике с Райлом) и на тарсковское понимание истины как соответствия фактам. Из этого берётся: (1) различение формальной истинности и эмпирического применения, (2) тезис об опровержимости применений, (3) роль логики как условия корректного следования, а не источника фактуального содержания.

\subsection{Контрольные точки с подробными ответами}

\begin{enumerate}
    \item В чём специфика формальных наук у Поппера?\\
    Они не описывают факты сами по себе; описательными и опровержимыми становятся их применения к реальности (через интерпретацию и эмпирические предпосылки).

    \item Почему $2+2=4$ «меняет статус» при применении?\\
    Внутри формальной арифметики это логическая истина, зависящая от определений и правил. В мире вещей это работает только при дополнительных допущениях о том, что именно мы считаем объектами и что с ними происходит; эти допущения могут не выполниться, поэтому применение может быть опровергнуто.

    \item Какова роль логики в этой картине?\\
    Логика не добавляет фактов о мире; она задаёт условия корректного следования (сохранение истинности вывода). Поэтому она служит инструментом критики и проверки наших утверждений о мире, а не «теорией фактов».
\end{enumerate}


\newpage
\section{Билет №5. Проблема критерия истины. Истина и знание}

\subsection{Ответ}
Проблема критерия истины состоит не просто в выборе «способа обоснования», а в необходимости найти такой признак (формальный или эмпирический), который позволял бы субъекту \emph{практически} отличать истину от заблуждения в процессе познания. Анкин указывает на критическую важность различения \emph{определения} истины и её \emph{критерия}: определение сообщает, что такое истина (например, соответствие реальности), тогда как критерий — это метод или тест её распознавания; смешение этих уровней приводит к ошибке, когда саму истину подменяют условиями её обнаружения или проверки.

Традиционно знание определяется как обоснованное истинное убеждение. Это означает, что для обладания знанием субъект должен не просто иметь в сознании истинную пропозицию, но и располагать достаточными доводами в её пользу. Проблема критерия здесь проявляется в вопросе о том, что именно считается достаточным обоснованием. Один из подходов предлагает критерий когерентности, согласно которому истинность тезиса подтверждается его логической непротиворечивостью относительно всей системы уже принятых убеждений. Другой подход — верификационизм, предполагающий, что критерием является возможность опытной проверки. Бертран Рассел критикует сведение истины к верификации, утверждая, что факт остается фактом даже в том случае, если у нас нет инструментов для его обнаружения.

Однако даже когда у субъекта есть «обоснование» и убеждение оказывается истинным, это ещё не гарантирует знания: именно это показывает проблема Гетье, где истинность получается за счёт случайного совпадения факторов (эпистемического везения). Поэтому вопрос о критерии обоснованности становится центральным: важно не просто иметь доводы, а понимать, \emph{насколько они исключают ошибку}.

Именно в этом смысле «критерий обоснованности остаётся относительным по отношению к вычислительным мощностям и опыту конкретного субъекта». Истина как соответствие миру не зависит от субъекта, но \emph{обоснование} зависит от того, какие альтернативы субъект вообще способен рассмотреть и исключить. В терминах Хинтикки знание требует исключения релевантных альтернатив тезису $p$, однако набор таких альтернатив и доступные способы их проверки задаются: \begin{enumerate}
    \item опытом и информацией субъекта (какие данные и наблюдения ему реально доступны);
    \item его познавательными ресурсами (насколько сложные рассуждения и какой объём информации он способен обработать).
\end{enumerate}
Поэтому один и тот же тезис может быть истинным объективно, но «достаточность» обоснования для признания его знанием будет различаться: у одного субъекта альтернативы исключены, у другого — нет, и именно здесь возникает пространство для гетьеровских случаев эпистемического везения.

\subsection{Обзор источников}
Бертран Рассел в главе Истина и верификация последовательно критикует прагматистские и верификационистские попытки заменить истину критерием успешности или проверяемости. Его аргументация строится на защите реальности: если истина — это только то, что верифицируемо, то мы теряем огромное количество фактов о прошлом. Яакко Хинтикка в работе Проблема истины анализирует критерии через призму семантики возможных миров. Он доказывает, что субъект обладает знанием только тогда, когда он может исключить все релевантные альтернативы исходному тезису. Анкин в Теории познания посвящает отдельные разделы разбору дистинкции между определением и критерием, а также детальному анализу проблемы Гетье. Он показывает, как попытки найти универсальный критерий истины в истории философии разбивались о логический скептицизм и регресс обоснования.

\subsection{Контрольные точки с подробными ответами}

\textbf{В чем принципиальное различие между определением истины и критерием истины?}
Определение задает смысл понятия истина как соответствия реальности, а критерий — это метод или тест, позволяющий практически разделить утверждения на истинные и ложные.
\textbf{Что такое эпистемическое везение в концепции Рассела?}
Это ситуация, когда убеждение познающего субъекта оказывается верным по случайному стечению обстоятельств, а не в силу правильного метода познания или надежного обоснования.
\textbf{Какова трехчастная структура определения знания в аналитической традиции?}
Знание включает в себя три обязательных компонента: наличие убеждения у субъекта (S верит, что p), объективную истинность этого убеждения (p истинно) и наличие веских оснований.
\textbf{Почему непротиворечивость часто считается недостаточным критерием истины?}
Согласно Расселу, можно построить систему высказываний, которая будет идеально согласованной (когерентной), как хороший роман, но при этом будет полностью вымышленной.


\newpage
\section{Билет №6. Понятие онтологических обязательств и другие базовые идеи в работе У. Куайна «О том, что есть»}

\subsection{Ответ}

Куайн в статье «О том, что есть» предлагает разбирать онтологические споры не на уровне интуиций («кажется, что есть такие-то сущности»), а через логический анализ того, \emph{что именно утверждает} наша теория в ясной (канонической) логической записи. Отправная точка --- платоновская трудность: когда мы говорим «Пегаса не существует», создаётся впечатление, что мы уже каким-то образом обязаны признать Пегаса (иначе непонятно, о чём речь). Куайн считает, что эта ловушка возникает из неверного допущения: будто всякое осмысленное выражение обязано быть именем некоторого объекта.

Чтобы снять проблему, Куайн опирается на расселовский приём анализа дескрипций: выражения, которые выглядят как имена, можно переписать так, чтобы они не требовали особого «предмета» в онтологии. В терминах логики это означает: вместо того чтобы обращаться с «Пегасом» как с именем (которое должно что-то обозначать), мы переводим высказывание в кванторную форму, где видно, что речь идёт о том, есть ли объект, удовлетворяющий определённому условию. Тогда фраза «Пегаса не существует» в итоге становится отрицанием утверждения вида $\exists x\, P(x)$ (где $P(x)$ --- условие вроде «$x$ есть пегас»), и нам не нужно вводить «несуществующие сущности», чтобы объяснить осмысленность отрицания.

Однако новизна Куайна по сравнению с Расселом не сводится к повторению теории дескрипций. Рассел показывает, \emph{как} устранять некоторые вводящие в заблуждение выражения; Куайн делает из этого общий метод онтологии: прежде чем решать, какие сущности признавать, надо регламентировать язык теории и посмотреть, какие сущности \emph{обязательны} для истинности её утверждений. То есть у Куайна расселовский анализ работает как часть программы: онтология читается с логической формы теории.

Эта программа фиксируется в его центральном понятии --- \textbf{онтологического обязательства}. Онтологические обязательства теории --- это те объекты, которые должны входить в область значений переменных, чтобы утверждения теории могли быть истинными. Отсюда и критерий Куайна: \emph{существовать --- значит быть значением переменной}. Важный смысл критерия в том, что он различает два уровня:
\begin{enumerate}
    \item \textbf{упоминание/употребление выражения} (мы можем понимать слово и пользоваться им в речи);
    \item \textbf{признание существования в теории} (мы обязуемся существованием тех сущностей, по которым квантифицируем: когда в канонической записи требуется $\exists x$ так, чтобы некоторые $x$ делали утверждение истинным).
\end{enumerate}
Поэтому мы не «несём ответственность» за всё, что фигурирует в обычном языке: обязательство появляется не из факта, что слово значимо, а из того, что в логической форме теории требуется объект как значение переменной.

Куайн иллюстрирует это критикой оппонентов (МакИкс и Вимен), которые считают, что раз выражения осмысленны, то должны быть соответствующие сущности (например, «возможные» объекты или универсалии). Куайн разводит значение и именование: можно понимать употребление выражения и при этом не допускать, что оно что-то называет. Именно поэтому спор о том, «есть ли» те или иные сущности (например, универсалии), он переводит в вопрос о том, \emph{какие кванторы и по чему} необходимы в лучшей формулировке нашей теории. Если теория в канонической записи вынуждена квантифицировать по некоторому типу сущностей, то она принимает соответствующие онтологические обязательства; если можно переформулировать её так, чтобы обойтись без них, то Куайн предлагает предпочесть более экономную схему (в духе бритвы Оккама) при сохранении объяснительной силы.

\subsection{Обзор источников}

Основным источником является статья У. Куайна О том, что есть. Работа структурирована как последовательное опровержение аргументов в пользу перенаселенной онтологии. Сначала автор вводит персонажей МакИкса и Вимена для представления типичных ошибок — признания бытия объектов в качестве идей или нереализованных возможностей. Затем следует изложение метода Рассела как средства устранения имен. Центральная часть текста посвящена обоснованию критерия онтологического обязательства через переменные квантификации. В финальной части Куайн анализирует спор об универсалиях, проводя параллели между средневековыми позициями (реализм, концептуализм, номинализм) и современными течениями в философии математики (логицизм, интуиционизм, формализм), завершая работу призывом к прагматическому выбору концептуальной схемы.

\subsection{Контрольные точки с подробными ответами}

\begin{enumerate}
    \item \textbf{В чем заключается ловушка Платона в контексте онтологии?}
    Она состоит в убеждении, что для того, чтобы отрицать существование чего-либо (например, Пегаса), нужно сначала признать, что это нечто в каком-то смысле есть, иначе отрицание будет лишено объекта. Куайн решает это через замену имен дескрипциями.
    \item \textbf{Как звучит критерий онтологического обязательства по Куайну?}
    Быть — значит быть значением переменной. Теория налагает онтологические обязательства в том случае, если она утверждает существование сущностей как значений переменных квантификации $\exists x$.
    \item \textbf{Как Куайн предлагает поступать с именами собственными?}
    Он предлагает превращать их в предикаты (например, Пегас становится предикатом пегасировать) и использовать теорию дескрипций, чтобы полностью исключить имена из логического анализа существования.
    \item \textbf{Как Куайн разграничивает значение и именование?}
    Значение — это лингвистическая характеристика слова (понимание смысла), а именование — это отношение слова к конкретному объекту в мире. Можно иметь значимое слово, которое ничего не именует.
    \item \textbf{Чем руководствуется Куайн при выборе между разными онтологиями?}
    Прагматическим критерием простоты и эффективности. Он сравнивает это с выбором научной гипотезы: мы выбираем ту систему объектов, которая лучше всего упорядочивает наш чувственный опыт.
\end{enumerate}


\newpage
\section{Билет №7. Проблема Э. Гетье и её влияние на понимание знания и обоснования}

\subsection{Ответ}

Проблема Гетье возникла после публикации в 1963 году короткой статьи Эдмунда Гетье «Является ли знанием обоснованное истинное убеждение?». В аналитической эпистемологии к этому моменту был широко принят классический анализ знания как обоснованного истинного убеждения ($JTB$), традиционно восходящий к платоновской линии рассуждений: чтобы субъект $S$ знал, что $p$, должно выполняться, что $S$ убеждён в $p$, $p$ истинно и у $S$ есть обоснование $p$. Гетье (в изложении у Гаспарова) показывает, что этих трёх условий недостаточно: можно получить ситуацию, где все они выполнены, но знанием это всё равно не является.

Суть контрпримеров Гетье состоит в том, что субъект получает истинное убеждение «слишком случайно»: истина достигается не благодаря качеству обоснования, а как удачное совпадение. Структурно (как это фиксируется у Гаспарова) часто используется схема, где у субъекта есть обоснованное убеждение $p$, при этом $p$ ложно, из $p$ логически следует $q$, и субъект обоснованно принимает $q$, но $q$ оказывается истинным по независимой причине. Классический пример (кейс Смита и Джонса): Смит имеет хорошие основания считать, что Джонс получит работу, и что у Джонса в кармане десять монет. Он делает вывод $q$: «человек, который получит работу, имеет десять монет в кармане». В действительности работу получает сам Смит, и у него (случайно) тоже десять монет. В итоге $q$ истинно, Смит в него верит и имеет обоснование (как ему кажется), но интуитивно он этого не \emph{знает}, потому что истинность вывода держится на удаче и опирается на ложное ключевое звено.

Влияние проблемы Гетье на эпистемологию связано с тем, что она выявила разрыв между логической корректностью/обоснованностью рассуждения и тем, что можно назвать «эпистемическим успехом без удачи» (эту линию подчёркивает Анкин). В ответ на гетьеровские контрпримеры (как обсуждается у Лебедева и Черняка, а также у Гаспарова) оформились две основные стратегии «спасения» анализа знания:
\begin{enumerate}
    \item \textbf{Добавить к $JTB$ четвёртое условие}, которое блокирует гетьеровские случаи. Сюда относят запрет на «ложные посылки» ($no\ false\ lemmas$): если в обосновании есть ложное звено, знание не получается. Однако даже если исключить явные ложные посылки, остаются случаи, где убеждение истинно и формируется стандартным способом, но всё равно истинность слишком случайна.
    \item \textbf{Пересмотреть само понимание обоснования/знания}, смещая акцент с внутренних доводов субъекта на свойства процесса формирования убеждения и условий среды (экстерналистские подходы). Релайабилизм формулирует это как требование: знание --- истинное убеждение, полученное надёжным когнитивным процессом. Тогда гетьеровские случаи трактуются как провалы надёжности: в специфическом контексте механизм, который привёл к убеждению, не обеспечивает устойчивого попадания в истину.
\end{enumerate}

Именно для демонстрации такого «контекстного» элемента удачи часто приводят кейс с «фальшивыми амбарами» (у Голдмана): субъект находится в местности, где почти все видимые «амбары» --- бутафорские фасады, не отличимые на глаз от настоящих, и лишь один амбар настоящий. Субъект смотрит как раз на настоящий амбар и формирует истинное убеждение $p$: «передо мной амбар» на основе обычного зрительного восприятия (без вывода из ложной посылки). Тем не менее интуитивно это всё ещё не знание: истинность получилась удачно, потому что в такой среде тот же способ формирования убеждений почти наверняка дал бы ошибку, если бы взгляд упал на очередной фасад. Этот пример показывает, что одной лишь корректной «внутренней» обоснованности может быть недостаточно: важны надёжность способа получения убеждения и отсутствие эпистемической удачи.

Таким образом, Гетье показал, что классическое обоснование в $JTB$-смысле не гарантирует устранения случайности: можно быть рационально обоснованным и прийти к истинному выводу, но прийти к нему «не тем путём», то есть без знания. Это и запускает «постгетьеровскую» перестройку теорий знания: либо вводятся дополнительные анти-гетьеровские условия, либо знание связывается с надёжностью когнитивных процессов и характеристиками эпистемической среды.

\subsection{Обзор источников}

Основной анализ проблемы содержится в разделе 9.2.1 работы \textbf{Гаспарова И. Г. «Knowledge»}. Он дает структурную схему аргумента Гетье: $p$ влечет $q$, $p$ обосновано, но ложно, $q$ случайно истинно. В источнике \textbf{Лебедева и Черняка} (глава «Аналитическая эпистемология») подробно разбираются стратегии выхода из кризиса: теории каузальной связи и теории «недефектного» обоснования. \textbf{Анкин Д. В.} (раздел 4.1) рассматривает проблему Гетье как логический парадокс, подчеркивая, что она выявила разрыв между логической обоснованностью и эпистемической удачей.

\subsection{Контрольные точки с подробными ответами}

\textbf{1. В чем главная логическая «дыра» JTB-модели, вскрытая Гетье?}\\
Модель допускает обоснованность ложных убеждений. Если из такого обоснованного, но ложного убеждения логически выводится истина, она формально считается знанием, хотя является случайной.

\textbf{2. Что такое условие «отсутствия ложных посылок»?}\\
Это предложенное дополнение к JTB, согласно которому знание требует, чтобы обоснование не опиралось ни на какие ложные допущения в ходе рассуждения.

\textbf{3. Как кейс с фальшивыми амбарами усложняет проблему Гетье?}\\
Он показывает, что знание может отсутствовать даже при истинных посылках и правильном восприятии, если субъект находится в среде, где вероятность ошибки слишком высока (эпистемическая среда ненадежна).

\textbf{4. К чему привел «постгеттиеровский» период в эпистемологии?}\\
К отказу от исключительного интернализма и расцвету экстерналистских концепций, таких как каузальная теория познания и когнитивная психология познания.


\newpage
\section{Билет №8. Основные категории теории познания И. Канта. Кант о границах познания.}

\subsection{Ответ}

Кант во «Введении» к работе «Критика чистого разума» ставит задачу выяснить, как возможно достоверное знание. Сначала он различает всё наше познание на априорное и апостериорное. Априорное знание мы имеем до и независимо от конкретного опыта; его признаки — необходимость и строгая всеобщность. Например, положение о том, что у каждого изменения есть причина, не выводится из простого наблюдения: опыт показывает лишь то, что бывает, но не то, что должно быть всегда. Апостериорное знание, наоборот, целиком опирается на чувственный опыт и поэтому носит случайный, ограниченный характер.

Далее Кант вводит различие аналитических и синтетических суждений. Аналитические суждения поясняют понятие: предикат уже содержится в понятии субъекта (например, «все тела протяженны»). Синтетические суждения расширяют знание: предикат добавляет новое. Ключевой ход Канта — утверждение существования синтетических априорных суждений: они одновременно расширяют знание и не зависят от опыта. Кант связывает с ними основания математики и теоретического естествознания. Например, в равенстве $5+7=12$ результат $12$ не «извлекается» простым анализом понятий $5$ и $7$; он получается через акт сложения, то есть через синтез, который (в кантовском объяснении) опирается на чистое созерцание времени как форму последовательного счёта. Поэтому это знание априорно, но синтетично.

Чтобы объяснить устройство познания, Кант различает две способности: чувственность и рассудок. Чувственность даёт нам предметы через созерцание; у неё есть две априорные формы (именно формы чувственности) — пространство и время. Важно: форм именно две, потому что Кант различает внешнее созерцание (пространство как форма внешних явлений) и внутреннее (время как форма внутреннего чувства, и в итоге как форма всякого явления для нас). Это не свойства «вещей самих по себе», а условия, при которых нам вообще может что-либо являться. Рассудок же мыслит данные созерцания с помощью априорных понятий — категорий (например, причинности, субстанции, единства): они выступают «правилами» синтеза, связывающими многообразие чувственного материала в предмет опыта. Поэтому рассудок относится к априорному уровню: категории не берутся из опыта, а являются условиями возможности опыта как упорядоченного знания.

Почему у Канта именно $12$ категорий и нужно ли их все перечислять: в самом замысле «Критики» Кант стремится получить список категорий не произвольно, а систематически — по числу и виду логических функций суждения (отсюда четыре группы и в каждой по три категории: количество, качество, отношение, модальность). Для ответа на билет обычно достаточно понимать \emph{принцип}: категории — не «список на удачу», а набор априорных понятий рассудка, организованный по логической форме суждений; знать все названия полезно, но по смыслу важнее уметь объяснить, что категории делают (делают возможным опыт как предметное знание) и привести 2--3 примера (причинность, субстанция и т.п.).

Границы познания у Канта определяются тем, что знание возможно только там, где соединяются чувственность и рассудок: мы познаём вещи лишь как они нам являются в формах пространства и времени и под категориями рассудка. Всё, что выходит за пределы возможного опыта, не может стать предметом теоретического знания. Для обозначения того, что мыслится «по ту сторону» явлений, Кант использует понятие ноумена (вещи самой по себе): оно служит указанием на предел, а не объектом, который можно познать. Попытки теоретического разума познать сверхопытное (например, Бога, душу, мир как целое) приводят к неразрешимым противоречиям, поэтому Кант ограничивает знание сферой возможного опыта, а вопросы сверхопытного переносит в область практического разума (морали).

\subsection{Контрольные точки с подробными ответами}

\begin{enumerate}
    \item \textbf{Почему арифметика для Канта является синтетической при всей её очевидности?}\\
    Потому что результат суммы не содержится аналитически в понятиях слагаемых: он получается через синтез (операцию сложения), который Кант связывает с чистым созерцанием времени как последовательностью прибавления единиц.
    \item \textbf{У созерцания только две априорные формы?}\\
    Да: у чувственности как способности созерцания две априорные формы — пространство и время. Это условия всякого возможного опыта для нас, а не свойства вещей самих по себе.
    \item \textbf{В чем заключается функция категорий рассудка и априорны ли они?}\\
    Категории — априорные понятия рассудка, с помощью которых он связывает многообразие чувственного материала в предмет опыта (например, через причинность или субстанцию). Поэтому рассудок участвует в априорном знании как источник этих форм мышления.
    \item \textbf{Почему категорий $12$ и важно ли знать их все по именам?}\\
    Кант стремится вывести категории систематически из логических функций суждения, поэтому получает четыре группы по три (количество, качество, отношение, модальность). На экзамене обычно важнее понимать принцип и роль категорий и уметь назвать несколько ключевых, чем механически перечислить все $12$.
    \item \textbf{Есть ли у Канта созерцание в аналитических априорных суждениях?}\\
    Аналитические суждения не расширяют знание и основаны на разъяснении понятий, поэтому они не требуют синтеза через (чистое) созерцание так, как это требуется в синтетических априорных суждениях (особенно в математике). Но знание о предметах опыта в целом у Канта возможно только при совместной работе чувственности (созерцания) и рассудка (категорий).
    \item \textbf{Может ли человек познать ноумен?}\\
    Нет: теоретическое познание ограничено условиями возможного опыта (пространством, временем и категориями). Ноумен указывает на предел этого познания и не дан в опыте.
\end{enumerate}


\newpage
\section{Билет №9. Категории априорного и апостериорного: от Канта до Крипке.}

\subsection{Ответ}

Категории априорного и апостериорного изначально вводятся как различение двух источников/способов обоснования знания. У Канта во «Введении» к «Критике чистого разума» априорным называется знание, не зависящее от опыта: оно предшествует всякому конкретному восприятию и обладает признаками необходимости и строгой всеобщности. Поэтому Кант противопоставляет его апостериорному (эмпирическому) знанию, которое целиком черпается из опыта и всегда носит случайный, неабсолютный характер. Типично: «трава зелёная» мы узнаём только наблюдением, а не из одного разума.

При этом у Канта важно не смешивать два различения: (i) априорное/апостериорное как характеристика \emph{нашего познания} и (ii) необходимость/случайность как характеристика \emph{истинности высказывания}. Во «Введении» Кант тесно связывает априорность с необходимостью и всеобщностью, поэтому в классической традиции часто складывается впечатление, что априорное $\Rightarrow$ необходимое, а апостериорное $\Rightarrow$ случайное.

В XX веке (что фиксируется у Лебедева в истории аналитической эпистемологии) эта связка подвергается пересмотру. Крипке в тексте «Тождество и необходимость» подчёркивает, что априорность и апостериорность --- это эпистемологические характеристики (как мы узнаём/обосновываем), а необходимость и случайность --- метафизические/модальные характеристики (могло ли быть иначе). Ошибка многих философских обсуждений, по Крипке, в том, что эти оси смешивали.

Ключевой ход Крипке --- показать, что возможны комбинации, которые кажутся «запрещёнными» при кантовском отождествлении априорного с необходимым.

\begin{enumerate}
    \item \textbf{Необходимые апостериорные истины.} Некоторые тождества мы узнаём только эмпирически, но если они истинны, то они истинны необходимо. Крипке иллюстрирует это через принцип необходимости тождества и через различие между именами/жёсткими обозначениями и описаниями. Классический пример: «Геспер = Фосфор». То, что Вечерняя звезда и Утренняя звезда --- одна и та же планета (Венера), устанавливается астрономически, то есть апостериорно. Но раз речь идёт об одном и том же объекте, то тождество объекта с самим собой не может быть иначе: оно необходимо. Аналогично устроены и случаи с естественными видами (в духе тождества «вода = $H_2O$»): знание об этом добывается опытно, но если термин фиксирует один и тот же природный вид, то тождество имеет модальный статус необходимости.
    \item \textbf{Случайные априорные истины.} Крипке также показывает, что можно знать некоторое утверждение априори, хотя оно не является необходимым. Типичный пример --- введение единицы измерения через «эталон»: если «$1$ метр» фиксируют как длину конкретного стержня в момент $t_0$, то утверждение «длина этого стержня в момент $t_0$ равна $1$ метру» известно априори (в силу правила фиксации значения), но оно случайно: стержень мог бы иметь другую длину, и тогда при той же процедуре фиксации «метр» был бы иным.
\end{enumerate}

Для обоснования этих выводов Крипке вводит различие \textbf{жёстких} и \textbf{нежёстких} обозначений. Жёсткий десигнатор (обычно собственное имя, а также многие термины естественных видов в его трактовке) обозначает один и тот же объект во всех контрфактических ситуациях, где этот объект существует; нежёсткие выражения (особенно дескрипции вроде «нынешний президент») могут указывать на разных людей в разных возможных ситуациях. Именно из жёсткости имён и принципа необходимости тождества Крипке получает результат: некоторые тождества познаются апостериорно, но истинны необходимо.

Итог «от Канта до Крипке» таков: кантовское различение априорного и апостериорного остаётся различением \emph{источника/способа знания}, но Крипке разводит его с модальными категориями. Теперь нельзя автоматически заключать, что априорное обязательно необходимо, а апостериорное обязательно случайно: это разные измерения, и между ними возможны перекрёстные сочетания.

\subsection{Контрольные точки с подробными ответами}

\begin{enumerate}
    \item \textbf{В чем основное различие между эпистемологией и метафизикой по Крипке?}\\
    Эпистемология отвечает, как мы узнаём и обосновываем истину (априори/апостериори). Метафизика отвечает, могла ли истина быть иной (необходимость/случайность).
    \item \textbf{Почему утверждение «Геспер есть Фосфор» метафизически необходимо, хотя познаётся апостериорно?}\\
    Потому что это тождество одного и того же объекта, установленное эмпирически; но тождество объекта с самим собой (при жёстком обозначении) не может быть иначе, то есть необходимо.
    \item \textbf{Что такое жёсткий десигнатор?}\\
    Это выражение, которое обозначает один и тот же объект во всех возможных ситуациях, где этот объект существует (в противоположность дескрипциям, которые могут «переключаться» на разные объекты).
    \item \textbf{Почему «эталон метра» даёт пример случайного априорного?}\\
    Потому что истинность «стержень в момент $t_0$ имеет длину $1$ метр» обеспечивается правилом фиксации значения и потому известна априори, но сама длина стержня могла бы быть другой, значит утверждение не необходимо.
\end{enumerate}

\newpage
\section{Билет №10. Категории аналитического и синтетического: от Канта до Куайна.}

\subsection{Ответ}

\subsubsection*{1. Кант: что такое аналитическое и синтетическое}

Во «Введении» к «Критике чистого разума» Кант различает два типа суждений:
\begin{enumerate}
    \item \textbf{Аналитические} — предикат уже содержится в понятии субъекта; такие суждения \emph{поясняют} понятие и опираются на закон противоречия. Пример: «треугольник имеет три угла».
    \item \textbf{Синтетические} — предикат добавляет к понятию субъекта новое; такие суждения \emph{расширяют} знание. Поэтому синтез Кант считает ключевой операцией научного знания (как подчёркивает Лебедев).
\end{enumerate}
У Канта это различение важно для общей задачи: объяснить возможность синтетических (в том числе синтетических априорных) оснований науки.

\subsubsection*{2. Переход к аналитической философии: запрос на строгий критерий}

В XX веке (в частности, у логических позитивистов) появляется стремление сделать различие строгим в терминах языка: аналитическое начинают понимать как «истинное в силу значения/правил употребления», а синтетическое — как «зависящее от фактов». Куайн в «Двух догмах эмпиризма» атакует именно эту идею чёткой границы, считая её философской «догмой».

\subsubsection*{3. Куайн: почему определения аналитичности круговые}

Куайн показывает, что стандартные объяснения аналитичности неизбежно замыкаются в круг. Упрощённо логика такая:
\begin{enumerate}
    \item Мы хотим определить \textbf{аналитичность} как истинность «в силу значения», потому что это попытка переформулировать кантовскую идею «предикат содержится в понятии субъекта» в языковых терминах: аналитическое должно быть истинным благодаря правилам употребления/значения слов и не требовать обращения к фактам опыта.
    \item Тогда нужно объяснить \textbf{значение} через \textbf{синонимию}: когда выражения равнозначны (например, «холостяк» и «неженатый мужчина»).
    \item Но когда мы пытаемся обосновать синонимию, мы обычно апеллируем к тому, что предложения типа «все холостяки неженаты» истинны \emph{в силу значения}, то есть \emph{аналитичны}.
\end{enumerate}
Получается круг: аналитичность объясняют через синонимию, а синонимию — через аналитичность. То же относится и к ссылкам на «определения»: определение фиксирует употребление, но само не даёт независимого критерия, почему именно такая фиксация является «истинной по смыслу», не предполагая уже различия аналитического и синтетического.

\subsubsection*{4. Холизм и тезис «нет чисто аналитического»}

Далее Куайн усиливает вывод общей картиной знания — \textbf{холизмом}: наука (и вообще система наших убеждений) работает как целая сеть. Опыт «проверяет» не отдельные предложения изолированно, а сеть в целом: при конфликте с опытом мы можем перераспределять, что удерживать, а что пересматривать, и это касается даже очень центральных утверждений (Лебедев интерпретирует это как радикальную линию пересматриваемости).

Отсюда становится понятнее тезис: «невозможно выделить чистые аналитические истины». У Куайна это означает:
\begin{enumerate}
    \item \textbf{нет строгого, некругового критерия}, который бы отделял класс истин, зависящих \emph{только} от языка, от истин, зависящих от мира;
    \item различие «аналитическое/синтетическое» (если им вообще пользоваться) становится \textbf{прагматическим и степенным}: какие утверждения мы фактически держим фиксированными и какой ценой готовы их пересматривать внутри общей теории.
\end{enumerate}

\subsubsection*{5. Итог: от Канта к Куайну}

\begin{enumerate}
    \item У Канта различие аналитического и синтетического встроено в проект обоснования науки и возможности расширяющего знания.
    \item У Куайна (в «Двух догмах») под сомнение ставится сама идея чёткой границы: попытки определить аналитичность оказываются круговыми, а холистическая картина знания размывает представление о «чисто языковых» истинах.
    \item Практический вывод (как акцентирует Анкин): многие «структурные» элементы знания удерживаются не потому, что имеют привилегию аналитичности, а потому что они наиболее удобны и эффективны в рамках нашей лучшей теории.
\end{enumerate}


\newpage
\section{Билет №11. Основные идеи работы «Две догмы эмпиризма» У. Куайна}

\subsection{Ответ}

В статье «Две догмы эмпиризма» Куайн критикует программу логического позитивизма, утверждая, что классический эмпиризм держится на двух сомнительных предпосылках («догмах»), и предлагает альтернативную картину знания.

\subsubsection*{1. Первая догма: различие аналитического и синтетического}

Первая догма --- вера в чёткую границу между:
\begin{enumerate}
    \item \textbf{аналитическими} истинами, которые якобы истинны «в силу значения» или «по определению»;
    \item \textbf{синтетическими} истинами, которые истинны благодаря фактам.
\end{enumerate}

Куайн показывает, что \emph{строго} определить аналитичность не получается: все стандартные попытки оказываются круговыми.
\begin{enumerate}
    \item Если говорить, что аналитичность основана на \textbf{синонимии} (например, «холостяк» = «неженатый мужчина»), то нужно объяснить, что такое синонимия.
    \item Если объяснять синонимию через \textbf{определения} (словарные/явные), то определения лишь фиксируют уже предполагаемую равнозначность и не дают независимого основания.
    \item Если объяснять синонимию через \textbf{взаимозаменимость без изменения истинности} (salva veritate), то выясняется, что такая взаимозаменимость работает лишь при дополнительных предпосылках о том, какие контексты допустимы и какие логические принципы уже принимаются; то есть снова незаметно предполагается то, что хотели определить.
    \item Попытка Карнапа ввести \textbf{семантические правила} искусственного языка, по Куайну, не решает проблему: чтобы такие правила действительно фиксировали «аналитичность», нужно уже понимать, что считается «значением» и «синонимией», иначе мы просто переименовали трудность.
\end{enumerate}

Отсюда вывод Куайна: аналитичность не получает некругового объяснения, поэтому представление о принципиальном разрыве «истины по смыслу» и «истины по факту» не имеет надёжного основания.

\subsubsection*{2. Вторая догма: редукционизм (верификация поодиночке)}

Вторая догма --- редукционизм, то есть идея, что каждое отдельное высказывание имеет собственный эмпирический «перевод» и проверяется опытом изолированно. Куайн возражает, что опыт соотносится с системой знаний как с целым: конфликт с наблюдением устраняется перестройкой разных участков теории, а не «приговором» одному предложению.

Вместо этого он формулирует \textbf{эпистемологический холизм}: единицей эмпирической проверки является не отдельное предложение, а \emph{связный фрагмент теории, вплоть до науки в целом}. Мы сталкиваемся с опытом как система: наблюдение вступает в конфликт не с одним предложением напрямую, а с некоторым участком всей «сети» наших убеждений. Поэтому при столкновении с «неподходящим» опытом у нас обычно есть выбор, \emph{что именно} пересматривать: можно отказаться от конкретной гипотезы, добавить вспомогательное допущение, уточнить условия применения и т.д. В этом смысле, по Куайну, \emph{нет утверждений, абсолютно иммунных от пересмотра}: даже логика и математика находятся ближе к центру сети и пересматриваются крайне неохотно, но принципиальной неприкосновенности не имеют.

\subsubsection*{3. «Эмпиризм без догм»: прагматические критерии и постулаты}

После отказа от двух догм Куайн предлагает «эмпиризм без догм»: наши теории остаются ориентированными на опыт, но связь с опытом опосредована всей системой. Поэтому выбор между альтернативными реконструкциями теории становится частично \textbf{прагматическим}: мы оцениваем, какая система проще, удобнее, лучше организует опыт и успешнее работает в целом.

Физические объекты у Куайна выступают как \textbf{постулаты} теории: они не «выписаны» прямо из чистых ощущений, а вводятся как удобные сущности, с помощью которых мы упорядочиваем чувственные данные и предсказываем дальнейший опыт. Поэтому он и сравнивает их с «богами Гомера»: и те, и другие являются постулируемыми сущностями, но физические объекты принимаются потому, что они дают значительно более простую, связную и успешную систему объяснения опыта.

\subsection{Контрольные точки с подробными ответами}

\begin{enumerate}
    \item \textbf{Почему Куайн считает аналитичность неопределимой?}\\
    Потому что попытки определить её через синонимию, определения, взаимозаменимость или семантические правила приводят к кругу или к скрытым предпосылкам, которые уже предполагают различие аналитического и синтетического.
    \item \textbf{В чем суть куайновского холизма?}\\
    В том, что наши высказывания сталкиваются с опытом не по одному, а как часть целой системы: проверяется сеть убеждений, и при конфликте возможны разные способы пересмотра внутри этой сети.
    \item \textbf{Какова роль логики и математики в системе Куайна?}\\
    Они ближе к центру «сети» и пересматриваются с наибольшей неохотой, но принципиально не обладают абсолютной неприкосновенностью: в пределе пересмотр возможен как системное решение.
    \item \textbf{Как Куайн относится к физическим объектам?}\\
    Как к полезным постулатам теории: они не «даны» в чистом виде в опыте, но принимаются как удобные посредники для организации чувственных данных; поэтому уместно сравнение с мифологическими постулатами, хотя научные постулаты эффективнее.
    \item \textbf{Что остается от эмпиризма после критики Куайна?}\\
    Остаётся ориентация на чувственный опыт как на «трибунал», но без редукции предложений поодиночке: опыт оценивает теорию целиком, а внутренняя перестройка системы в значительной мере руководствуется общими (прагматическими) достоинствами теории.
\end{enumerate}


\newpage
\section{Билет №12. Категории смысла и значения (Г. Фреге)}

\subsection{Ответ}

Фреге в работе Смысл и значение приступает к анализу категории равенства, задаваясь вопросом о том, является ли оно отношением между предметами или между их знаками. Если бы равенство $a = b$ было просто отношением между предметами, то оно ничем не отличалось бы от $a = a$. Однако Фреге замечает, что суждение $a = a$ аналитично и дано априори, тогда как $a = b$ часто содержит новое знание, полученное опытным путем (апостериори). Это различие в познавательной ценности заставляет Фреге постулировать, что у знака, помимо обозначаемого предмета, есть еще одна характеристика — способ его данности.

Фреге вводит строгое различие между тремя уровнями: знаком, смыслом и значением. Значением имени является сам предмет, который оно обозначает. Смыслом же является способ данности этого предмета, тот интеллектуальный путь, который ведет нас к значению. Например, выражения утренняя звезда и вечерняя звезда имеют одно и то же значение — планету Венера, но их смыслы различны, так как они представляют этот объект в разных контекстах наблюдения. Фреге подчеркивает, что смысл является объективным: это не субъективное представление в уме индивида, а некое общее достояние, которое может быть понятно разным людям. Представление же — это психологический образ, уникальный для каждого человека и непередаваемый в языке.

Для Фреге идеальная связь заключается в том, что каждому знаку соответствует один смысл, а каждому смыслу — одно значение. Однако в естественных языках это правило часто нарушается. Например, имя Одиссей имеет смысл, но может не иметь значения, если Одиссей — лишь плод поэтического воображения. Фреге доказывает, что в искусстве (эпосе) нас волнует только смысл, создающий художественную реальность. Однако в науке мы требуем наличия значения, так как нас интересует истинностный статус высказывания.

Развивая свою теорию, Фреге переходит от имен к анализу целых предложений. Смыслом предложения является выраженная в нем мысль. Значением же предложения выступает его истинностное значение — Истина или Ложь. Куайн и Рассел позже будут критиковать этот ход, но для Фреге это логическое следствие принципа подстановочности: если мы заменим в предложении одно имя другим с тем же значением, но иным смыслом, мысль изменится, но истинность предложения останется прежней. Таким образом, истина — это тот абстрактный объект, на который указывает предложение. Особое внимание Фреге уделяет косвенному употреблению слов (в подчиненных предложениях), где значением слов становится их обычный смысл. Это объясняет, почему в контексте прямой или косвенной речи мы не можем просто заменять равнозначные термины без риска изменить истинность всего высказывания.

\subsection{Обзор источников}

Фрагмент работы Г. Фреге Смысл и значение начинается с анализа парадокса тождества на примере астрономических открытий. Текст структурно разделен на несколько этапов аргументации. Сначала автор вводит дистинкцию между смыслом и значением для имен собственных. Затем он проводит различие между объективным смыслом и субъективным представлением (используя знаменитую метафору с Луной и телескопом). Вторая половина статьи посвящена распространению этих категорий на повествовательные предложения. Фреге последовательно обосновывает, почему мысль — это смысл, а истина — это значение. Завершается фрагмент анализом сложных предложений и косвенной речи, где вводятся понятия прямого и косвенного значений.

\subsection{Контрольные точки с подробными ответами}

\begin{enumerate}
    \item \textbf{В чем разница между познавательной ценностью $a=a$ и $a=b$?}
    Суждение $a=a$ тривиально и истинно априори. Суждение $a=b$ (например, вечерняя звезда есть утренняя звезда) содержит эмпирическое открытие, так как сообщает, что два разных способа данности (смысла) относятся к одному и тому же предмету (значению).
    \item \textbf{Что такое смысл по Фреге и чем он отличается от представления?}
    Смысл — это объективное содержание мысли, доступное многим. Оно не зависит от воображения индивида. Представление — это субъективный чувственный образ в душе человека, который у каждого свой и не может быть частью логического анализа.
    \item \textbf{Как метафора с телескопом поясняет теорию Фреге?}
    В этой метафоре Луна — это значение (объект), реальный образ в линзе телескопа — это смысл (объективен для всех наблюдателей), а изображение на сетчатке глаза — субъективное представление.
    \item \textbf{Что Фреге называет значением предложения?}
    Значением предложения является его истинностное значение (Истина или Ложь). Это связано с тем, что нас интересуют факты, а не просто комбинации смыслов, когда мы ищем истину.
    \item \textbf{Что происходит со смыслом и значением в косвенной речи?}
    В косвенной речи происходит смещение: обычный смысл слова становится его значением (косвенным значением). Это необходимо для анализа таких контекстов, как Петр верит, что..., где истинность зависит от содержания мысли, а не от реальных объектов.
\end{enumerate}

\newpage
\section{Билет №13. Возможные миры, жёсткие десигнаторы и апостериорная необходимость (С. Крипке)}

\subsection{Ответ}

\subsubsection*{1. Крипке против дескриптивной теории имён}

В работе «Тождество и необходимость» Крипке критикует дескриптивную теорию имён, идущую от Фреге и Рассела. В рамках дескриптивизма имя рассматривается как «сокращённая дескрипция» или как выражение, смысл которого задаётся описанием: например, «Аристотель» понимается как что-то вроде «учитель Александра Македонского» (или набор таких описаний). Почему у дескриптивистов вообще возникает эта идея: они пытаются объяснить, за счёт чего имя \emph{что-то обозначает} и почему тождества с именами могут быть информативными. Если значение имени задаётся описанием, то понятно, почему «Аристотель = учитель Александра» может сообщать новое: оно связывает имя и дескрипцию. Именно против такого объяснения работы имён Крипке и выстраивает свои аргументы.

\subsubsection*{2. Возможные миры и контрфактические ситуации: что это и где границы}

Крипке вводит разговор о возможных мирах как разговор о контрфактических ситуациях: «что было бы, если бы было иначе». Это не «параллельная планета», а способ рассуждать о модальности: мы задаём условия и спрашиваем, могло ли при них быть так-то.

Граница здесь не в том, что мы можем \emph{вообразить} что угодно, а в том, что рассматриваются ситуации, которые считаются \emph{возможными} (в модальном смысле) при фиксированных правилах тождества и референции: когда мы обсуждаем «что могло бы случиться с данным объектом», мы должны удерживать, о каком именно объекте идёт речь. Поэтому контрфактическая ситуация — это не просто «какая-то история», а история, в которой мы продолжаем говорить о тех же предметах (если они в данной ситуации существуют) и спрашиваем, какие свойства/события могли бы быть иными.

\subsubsection*{3. Жёсткие и нежёсткие десигнаторы: почему это снимает путаницу}

Центральное понятие у Крипке — \textbf{жёсткий десигнатор}: термин, который обозначает один и тот же объект во всех возможных мирах, в которых этот объект существует. Крипке утверждает, что собственные имена (например, «Ричард Никсон») работают именно так.

В отличие от этого, дескрипции вроде «победитель выборов 1968 года» обычно \textbf{нежёсткие}: в другом возможном мире победителем мог оказаться другой человек, и тогда дескрипция будет указывать на другого.

Это напрямую отвечает на твою интуицию «в других воображаемых мирах мы же о разных людях говорим». Крипке говорит: зависит от того, \emph{каким выражением} ты пользуешься.
\begin{enumerate}
    \item Если ты говоришь: «Никсон мог проиграть», то при жёсткой референции речь идёт \emph{о том же Никсоне}, просто приписывается другой исход (проигрыш) в контрфактической ситуации.
    \item Если ты говоришь: «победитель выборов 1968 мог проиграть», то в другом мире «победитель» может быть другим человеком, и тогда утверждение уже относится к \emph{другому} человеку.
\end{enumerate}
То есть ощущение «мы говорим о другом объекте» обычно возникает тогда, когда незаметно подменяют имя (жёсткий десигнатор) дескрипцией (нежёсткий десигнатор).

Важно также уточнение, которое помогает отпустить последнюю мысль: в некоторых возможных мирах Никсон \emph{может не существовать} (Крипке это допускает как нормальный случай). Тогда имя не «переключается» на другого человека: оно просто не имеет там денотата. Жёсткость означает: \emph{если} объект существует в данном мире, имя продолжает обозначать именно его, а не «кого-то подходящего».

\subsubsection*{4. Апостериорная необходимость}

Из жёсткости имён и принципа необходимости тождества Крипке выводит тезис о \textbf{необходимых истинах, познаваемых апостериори}. Он разделяет:
\begin{enumerate}
    \item \textbf{априорность/апостериорность} как эпистемический статус (как мы узнаём истину);
    \item \textbf{необходимость/случайность} как модальный статус (могло ли быть иначе).
\end{enumerate}

Если $a$ и $b$ — жёсткие десигнаторы и тождество $a=b$ истинно, то оно необходимо истинно. Пример: «Вечерняя звезда = Утренняя звезда» (то есть обе — Венера). Мы узнаём это из опыта, значит истина апостериорна. Но раз речь об одном и том же объекте, при жёстком обозначении тождество не может «сломаться» в другом возможном мире: оно необходимо.

Крипке распространяет анализ и на термины естественных видов и научные тождества: знание о них может быть эмпирическим, но при установлении тождества (в рамках его подхода к референции) оно получает модальный статус необходимости. Иллюзия «могло бы быть иначе», по Крипке, часто возникает из смешения самого явления с нашим способом его обнаружения (то, что мы узнаём через опыт, не делает истину автоматически случайной).

\subsection{Контрольные точки с подробными ответами}

\begin{enumerate}
    \item \textbf{Почему дескриптивисты считали имя «краткой дескрипцией»?}\\
    Потому что они связывали значение имени с описанием (или набором описаний), которое, как предполагается, фиксирует референцию и объясняет информативность высказываний тождества с именами.
    \item \textbf{Как Крипке понимает возможные миры/контрфактические ситуации и где границы?}\\
    Как способы описания того, что могло бы быть иначе; граница в модальной возможности при сохранении корректной фиксации того, о каком объекте идёт речь (и без «переключения» имени на другого объект).
    \item \textbf{Что такое жёсткий десигнатор?}\\
    Термин, обозначающий один и тот же объект во всех возможных мирах, где этот объект существует; собственные имена у Крипке — парадигмальный пример.
    \item \textbf{Почему в других мирах мы не «переходим к другому человеку», говоря «Никсон мог проиграть»?}\\
    Потому что имя «Никсон» жёстко удерживает референцию: оно говорит об одном и том же человеке в разных контрфактических ситуациях; «другой человек» появляется при замене имени дескрипцией вроде «победитель выборов 1968 года».
    \item \textbf{В чем суть апостериорной необходимости?}\\
    Есть истины, которые узнаются только из опыта, но если они истинны, то истинны необходимо (например, тождества с жёсткими десигнаторами, такие как «Геспер = Фосфор»).
\end{enumerate}


\newpage
\section{Билет №14. Проблема существования в трудах Г. Фреге и Дж. Э. Мура}

\subsection{Ответ}

В аналитической философии вопрос «является ли существование предикатом?» означает: можно ли понимать слово «существует» как обычный признак вещи (как «рычит», «красный», «тяжёлый»), или его логическая роль иная.

\subsubsection*{1. Фреге: существование как предикат второго порядка}

Ключевой тезис Фреге состоит в том, что существование не является предикатом первого порядка (не приписывается отдельному объекту), а относится к понятиям, то есть является предикатом второго порядка. Когда мы говорим «Собаки существуют», мы тем самым говорим не о какой-то одной собаке и не приписываем собакам дополнительное свойство «существовать». Мы утверждаем, что понятие \emph{собака} не пусто: есть хотя бы один предмет, который под него подпадает.

В логической форме это выражается квантором существования:
$$\exists x\, Dog(x).$$
Здесь «существует» не выступает как предикат вида $Exist(x)$, а задаёт утверждение о том, что у предиката $Dog(x)$ есть хотя бы одно значение $x$, при котором он истинен. Поэтому, с точки зрения Фреге, формулировки вроде «существование — свойство предмета» вводят в заблуждение: логически корректнее говорить о непустоте понятия.

\subsubsection*{2. Мур: почему «существует» не похоже на обычные предикаты}

В статье «Является ли существование предикатом?» Мур показывает различие между «существует» и обычными приписываемыми свойствами на примерах.

Сравнение работает так:
\begin{enumerate}
    \item «Все дрессированные тигры рычат» — осмысленное обобщение: оно сообщает, какое свойство имеют тигры данного вида.
    \item «Все дрессированные тигры существуют» — звучит странно, потому что не добавляет к тиграм нового признака: если речь идёт о настоящих тиграх, то их существование уже подразумевается самим тем, что мы о них говорим как о тиграх в мире.
\end{enumerate}

Ещё яснее это видно на отрицаниях:
\begin{enumerate}
    \item «Некоторые тигры не рычат» — обычное, понятное отрицание: мы берём класс тигров и говорим, что не всем им присуще рычание.
    \item «Некоторые тигры не существуют» — в обычном смысле слова «тигр» это либо бессмыслица, либо требует смены контекста: чтобы фраза стала осмысленной, «тигры» должны пониматься не как реальные животные, а как, например, тигры, о которых мы говорим в воображении/рассказе/мифе.
\end{enumerate}

Именно отсюда Мур делает вывод: «существует» не работает как предикат «в одном ряду» с «рычит», потому что для обычных предметных терминов нельзя просто так образовать содержательное противопоставление «тигр, который существует» / «тигр, который не существует» (в отличие от «тигр, который рычит» / «тигр, который не рычит»).

\subsubsection*{3. Почему у Мура всё же бывают контексты, где «существует» похоже на предикат}

Мур подчёркивает, что слово «существует» действительно ведёт себя необычно, но при этом в некоторых контекстах мы употребляем его так, будто оно «отбирает» из одного множества элементы, которые есть в реальности, и элементы, которых в реальности нет.

Например, когда мы говорим о воображаемых объектах, становится осмысленным предложение вида:
\begin{quote}
«Некоторые из тигров, которых я себе вообразил, существуют (то есть есть такие тигры в реальности), а некоторые — нет».
\end{quote}
Здесь «существуют» не приписывает свойство тиграм как «вещам вообще», а выполняет роль различителя между двумя режимами: \emph{тигры как содержимое воображения} и \emph{тигры как реальные объекты}. Мур использует такие примеры, чтобы показать: некоторая «предикатность» слова «существует» появляется, когда мы заранее задали более широкий контекст (например, множество воображаемых объектов), внутри которого потом выделяем те, что имеют соответствие в действительности.

Отдельный важный для Мура случай — высказывания вида «Это существует», когда «это» указывает на непосредственно данное в опыте (sense-data, например, воспринимаемое пятно/цветовой образ). В таком контексте фраза становится осмысленной, потому что мы не рассуждаем о «тиграх вообще», а фиксируем наличие данного переживания: отрицание «этого нет» в момент его наличия выглядит как отрицание того, что непосредственно дано.

\subsubsection*{4. Отрицательные экзистенциальные высказывания и сближение с Фреге}

Мур также обсуждает парадоксальность отрицательных экзистенциальных предложений, вроде «Кентавры не существуют». Если понимать «существует» как обычный предикат, возникает ловушка: чтобы сказать, что у кентавров нет свойства существования, нужно было бы сначала иметь кентавров как носителей свойств. Поэтому (в русле расселовского анализа, к которому Мур обращается) такие предложения следует понимать не как «рассказ о кентаврах», а как утверждение о пустоте соответствующего понятия: нет объектов, подпадающих под «кентавр». Это сближает Мура с фрегевской линией: существование корректнее выражать через кванторную форму, а не как «свойство вещи».

\subsubsection*{Вывод}

Фреге даёт строгий логический ответ: существование — предикат второго порядка (непустота понятия), выражаемый квантором $\exists$. Мур подтверждает, что в обычной речи «существует» не функционирует как стандартный предикат, но уточняет, что в некоторых контекстах (воображаемые объекты, указание на непосредственно данное) употребление «существует» требует более тонкого анализа. В результате оба сходятся в критике идеи, что существование — обычное свойство объекта, и в объяснении отрицательных экзистенциальных высказываний через «нет объектов такого рода».

\subsection{Контрольные точки с подробными ответами}

\begin{enumerate}
    \item \textbf{Что означает тезис Фреге «существование — предикат второго порядка»?}\\
    «Существует» не приписывает свойство отдельной вещи, а утверждает непустоту понятия: $ \exists x\, F(x)$ означает, что есть хотя бы один предмет, подпадающий под $F$.
    \item \textbf{Как Мур показывает различие между «существует» и «рычит»?}\\
    Через тест на осмысленность обобщений и отрицаний: «некоторые тигры не рычат» понятно, а «некоторые тигры не существуют» в обычном смысле слова «тигр» не даёт нормального содержания и требует смены контекста (например, «тигры, которых я вообразил»).
    \item \textbf{В каком смысле, по Муру, фраза «Это существует» осмысленна?}\\
    Когда «это» указывает на непосредственно данное (sense-data): тогда высказывание фиксирует наличие данного переживания/объекта восприятия в момент опыта.
    \item \textbf{Как Фреге и Мур относятся к онтологическому доказательству бытия Бога?}\\
    Критически: если существование не является свойством объекта, его нельзя включать в перечень «совершенств» и выводить реальность из одного анализа понятия.
    \item \textbf{Зачем Мур обсуждает воображаемые объекты (например, единорогов/тигров)?}\\
    Чтобы показать контекст, в котором «существует» употребляется как различение: из множества воображаемых объектов выделяются те, которые имеют соответствие в реальности, и те, которые его не имеют.
\end{enumerate}

\newpage
\section{Билет №15. Семантическая концепция истины А. Тарского}

\subsection{Ответ}

Семантическая концепция истины А. Тарского направлена на то, чтобы дать строгое (логически корректное) определение истины для формализованных языков и тем самым устранить семантические парадоксы. Исходная трудность связана с \textbf{семантической замкнутостью} естественных языков: один и тот же язык способен
\begin{enumerate}
    \item описывать мир,
    \item говорить о собственных выражениях (называть предложения),
    \item приписывать им оценку «истинно/ложно».
\end{enumerate}
Это открывает путь самореференции и парадоксу лжеца.

\subsubsection*{1. Разделение языков: объектный язык и метаязык}

Решение Тарского состоит в строгом разграничении уровней:
\begin{enumerate}
    \item \textbf{объектный язык} $L$ --- язык, для которого мы хотим определить истину (в нём формулируются утверждения о некоторой предметной области);
    \item \textbf{метаязык} --- язык, в котором мы \emph{говорим о} выражениях $L$ и вводим предикат «истинно для $L$».
\end{enumerate}
Ключевое требование: определение истины для $L$ формулируется \emph{только} в метаязыке, чтобы избежать самоприменения предиката истины внутри одного и того же языка.

\subsubsection*{2. Т-схема (Конвенция $T$): условие материальной адекватности}

Тарский задаёт критерий правильности определения истины через \textbf{Т-схему}:
$$\ulcorner p\urcorner \text{ истинно } \Leftrightarrow p.$$
Здесь $p$ --- предложение объектного языка, а $\ulcorner p\urcorner$ --- \emph{имя} этого предложения в метаязыке (обычно запись предложения в кавычках). Например:
$$\text{«Снег бел» истинно } \Leftrightarrow \text{снег бел}.$$
Иногда в пересказах пишут «$X$ истинно тогда и только тогда, когда $p$»; это допустимо, если $X$ понимается именно как \emph{имя} предложения $p$. Смысл Т-схемы не в обозначениях буквами, а в различении ``имя предложения'' / ``само предложение''.

Т-схема задаёт не «механизм вычисления истины», а \textbf{требование}: адекватное определение истины должно порождать такие эквивалентности для всех предложений объектного языка.

\subsubsection*{3. Почему нужна интерпретация (модель)}

Чтобы вообще говорить об истинности формул объектного языка, Тарский рассматривает их \textbf{в интерпретации} (модели) $M$. Это не внешняя добавка, а часть семантики: нелогические символы языка сами по себе не фиксируют, какие объекты и отношения имеются в виду, пока не задано, \emph{как они интерпретируются}. Интерпретация включает:
\begin{enumerate}
    \item \textbf{домен} (предметную область) $D$;
    \item значения нелогических символов, в частности: каждому $n$-местному предикатному символу $R$ сопоставляется отношение $R^M \subseteq D^n$ (а константам и функциям, если они есть, --- соответствующие объекты и функции на $D$).
\end{enumerate}
После этого становится определимо, при каких подстановках переменных атомарные формулы истинны/выполнимы.

\subsubsection*{4. Выполнимость (удовлетворение) и рекурсивное определение истины}

Поскольку формул бесконечно много, Тарский определяет истину \emph{не перечислением}, а \textbf{рекурсивно по построению формулы}. Технической базой служит понятие \textbf{выполнимости} (удовлетворения): формула со свободными переменными выполнима в $M$ при заданном назначении переменных $s$ тогда, когда при этой подстановке она истинна.

Схема определения такова:
\begin{enumerate}
    \item Для \textbf{атомарных} формул (типа $R(x,y)$) выполнимость задаётся через интерпретацию предиката как отношения на $D$.
    \item Для \textbf{сложных} формул выполнимость определяется через выполнимость частей по правилам логических связок.
    \item Для \textbf{кванторов} правила задаются через домен $D$:
    \begin{enumerate}
        \item $\exists x\,\varphi(x)$ истинно в $M$ тогда и только тогда, когда существует $a \in D$ такое, что $\varphi(x)$ выполнима при назначении, где $x$ получает $a$;
        \item $\forall x\,\varphi(x)$ истинно в $M$ тогда и только тогда, когда для всех $a \in D$ выполняется соответствующая выполнимость $\varphi(x)$.
    \end{enumerate}
\end{enumerate}

Далее истина для \textbf{предложений} (формул без свободных переменных) получается как частный случай этой процедуры: для них уже не требуется подстановка свободных переменных, и можно говорить просто «истинно в интерпретации $M$».

\subsubsection*{5. Связь выполнимости с Т-схемой}

Связь такая: \textbf{Т-схема} задаёт критерий материальной адекватности («как должно получаться»), а \textbf{выполнимость} даёт способ построить строгое определение истины для формализованного языка (особенно с кванторами) так, чтобы необходимые экземпляры Т-схемы в метаязыке действительно выполнялись. Тем самым Тарский сохраняет классическую интуицию истины как соответствия, но формализует её для дедуктивных наук и избегает парадоксов за счёт разнесения уровней языка.


\newpage
\section{Билет №16. Работа Б. Рассела «Об обозначении» (основные идеи). Решение проблемы отрицательных экзистенциальных высказываний}

\subsection{Ответ}

В статье «Об обозначении» Рассел предлагает логический анализ тех выражений естественного языка, которые \emph{как будто} называют объект, хотя объект может отсутствовать. Его цель --- показать, как можно осмысленно говорить о «несуществующем» (например, о нынешнем короле Франции), не вводя в онтологию особые сущности и не нарушая законы логики. В этом контексте он полемизирует с Мейнонгом (у которого «объекты мысли» как бы есть даже при отсутствии существования) и с Фреге (различение смысла и значения), считая, что правильный выход даёт именно логическое разложение предложений.

\subsubsection*{1. Что такое обозначающие фразы}

\textbf{Обозначающие фразы} (denoting phrases) у Рассела --- это выражения, которые служат для «указания на предмет», но делают это не как собственные имена, а через общую форму вроде:
\begin{enumerate}
    \item \textbf{неопределённая} форма: «\emph{какой-то} человек», «\emph{некоторый} $x$»;
    \item \textbf{определённая} форма: «\emph{определённый} такой-то», «\emph{тот самый}», «\emph{нынешний король Франции}»;
    \item \textbf{универсальная} форма: «\emph{всякий} человек», «\emph{все} $x$, такие что \dots».
\end{enumerate}
Общее в них то, что они «обозначают» не прямым именованием, а через описание/кванторные слова, и именно поэтому могут порождать трудности, когда подходящего объекта нет или когда кажется, что выражение должно вести себя как имя.

\subsubsection*{2. Что такое дескрипции и почему они «неполные символы»}

Главный инструмент Рассела --- теория \textbf{дескрипций} (описаний). Дескрипции бывают двух типов:
\begin{enumerate}
    \item \textbf{неопределённые дескрипции}: «какой-то $F$» (в логике это соответствует $\exists x\,F(x)$);
    \item \textbf{определённые дескрипции}: «тот (единственный) $F$», например «нынешний король Франции», «автор \emph{Уэверли}».
\end{enumerate}

Ключевая идея Рассела: определённая дескрипция \emph{не является именем} и не имеет самостоятельного значения вне контекста; он называет её \textbf{неполным символом}. Это значит, что смысл выражения раскрывается только после того, как мы перепишем \emph{всё предложение}, где оно стоит, в логическую форму без «фиктивного имени».

Стандартное расселовское разложение предложения вида «\emph{определённый} $F$ есть $G$» имеет три компонента:
\begin{enumerate}
    \item \textbf{существование}: есть хотя бы один $x$ такой, что $F(x)$;
    \item \textbf{единственность}: есть не более одного $x$ с $F(x)$;
    \item \textbf{предикация}: этот $x$ имеет свойство $G$.
\end{enumerate}
Поэтому в примере «Автор \emph{Уэверли} был человеком» Рассел предлагает читать высказывание не как приписывание свойства «быть человеком» некоторому объекту-имени «автор \emph{Уэверли}», а как кванторное утверждение: существует ровно один автор \emph{Уэверли}, и он человек.

\subsubsection*{3. Решение проблемы отрицательных экзистенциальных высказываний}

Благодаря этому анализу Рассел решает проблему фраз вида «Нынешний король Франции не существует». Если бы «нынешний король Франции» был именем, мы попадали бы в тупик: чтобы отрицать существование объекта, пришлось бы как будто сначала иметь этот объект как значение имени. Но поскольку это \emph{дескрипция}, отрицательное экзистенциальное высказывание превращается в отрицание кванторной части: не существует такого $x$, который был бы нынешним королём Франции (или, в случае определённой дескрипции, не выполняются условия существования и единственности).

С этим связано расселовское различие \textbf{первичного} и \textbf{вторичного} вхождения дескрипции: в зависимости от того, на что распространяется отрицание, меняется логическая форма и, соответственно, истинностное значение. Это позволяет объяснить, почему некоторые предложения с несуществующими «объектами» оказываются не бессмысленными, а просто ложными или истинными в зависимости от области действия отрицания.

\subsubsection*{4. Непосредственное знакомство и «подлинные имена»}

В конце Рассел связывает теорию дескрипций с эпистемологией. Он различает:
\begin{enumerate}
    \item \textbf{непосредственное знакомство} (acquaintance) --- прямое знание того, что дано субъекту без посредства описаний (типичный пример у Рассела: данные опыта, на которые можно указать как «это»);
    \item \textbf{знание по описанию} (knowledge by description) --- знание объектов через описательные характеристики (например, «автор \emph{Уэверли}», «центр масс Солнца»).
\end{enumerate}
Рассел делает вывод, что многие слова, которые в обычной речи выглядят как имена собственные, в логическом анализе часто оказываются \emph{скрытыми дескрипциями}: они вводят объект не напрямую, а через набор описаний. Поэтому «подлинные» имена в строгом логическом смысле возможны лишь там, где референция опирается на непосредственное знакомство; всё остальное мы схватываем как «объект, удовлетворяющий описанию».

\subsection{Контрольные точки с подробными ответами}

\begin{enumerate}
    \item \textbf{Что такое обозначающие фразы у Рассела?}\\
    Это выражения с кванторными словами («некоторый», «всякий», «тот самый»), которые «указывают на предмет» через описание, а не как собственные имена.
    \item \textbf{Что такое дескрипции?}\\
    Это обозначающие фразы описательного типа; определённые дескрипции («тот единственный $F$») Рассел анализирует как неполные символы и устраняет через кванторы (существование, единственность, предикация).
    \item \textbf{Как теория дескрипций решает отрицательные экзистенциальные высказывания?}\\
    Она показывает, что отрицание относится не к «объекту имени», а к кванторному утверждению: в логической форме мы отрицаем существование объекта, удовлетворяющего описанию.
    \item \textbf{Что такое непосредственное знакомство и почему оно важно?}\\
    Это прямое знание того, что дано субъекту без описаний (на что можно указать как «это»); оно задаёт область «подлинных имён», тогда как большинство предметов познаётся по описанию.
\end{enumerate}


\newpage
\section{Билет №17. Проблема индукции по К. Попперу}

\subsection{Ответ}

Поппер рассматривает проблему индукции как центральную трудность классической эпистемологии и методологии науки. Традиционная картина (которую он связывает с индуктивизмом в духе Бэкона) предполагает, что наука движется от наблюдений единичных фактов к универсальным законам. Однако, следуя юмовскому аргументу, Поппер подчёркивает: \textbf{логического перехода} от конечного набора наблюдений к общему закону нет. Из того, что Солнце всходило каждое утро, не следует с логической необходимостью, что оно взойдёт завтра: всегда остаётся возможность исключения.

Отсюда Поппер делает вывод, что \textbf{индукция не может быть рационально обоснована}. Если попытаться оправдать индуктивный принцип через опыт (мол, «индукция раньше работала, значит, будет работать и дальше»), мы уже используем индуктивное рассуждение для обоснования индукции, то есть попадаем либо в круг, либо в бесконечный регресс «индукции более высокого порядка». Поэтому Поппер предлагает отказаться от идеи, что наука получает теории путём обобщения наблюдений. По его модели научное знание развивается иначе: мы выдвигаем \textbf{предположения} (гипотезы), а затем подвергаем их \textbf{критике} и попыткам опровержения.

Ключевой элемент этой альтернативы --- замена верификационизма критерием \textbf{фальсифицируемости}. Поппер указывает на асимметрию между подтверждением и опровержением универсальных утверждений. Никакое число совпадающих наблюдений не делает универсальный закон логически доказанным, но один противоречащий факт (при корректной фиксации) логически несовместим с законом. В простейшей форме это выражается через $modus\ tollens$: если из теории $T$ следует предсказание $P$, и обнаруживается $\neg P$, то $T$ отвергается.

Отсюда следует попперовское понимание научности: теория является научной не потому, что её «подтвердили», а потому, что она \textbf{заранее} допускает возможность своего опровержения. Фальсифицируемость означает наличие класса потенциальных фальсификаторов --- базисных высказываний наблюдения, которые, будучи приняты, противоречат теории. Поэтому развитие науки описывается как процесс устранения ошибок: теории остаются предположительными, даже если долго выдерживают проверки, и научная рациональность выражается не в накоплении подтверждений, а в систематической критике и отборе более устойчивых гипотез.

\subsection{Контрольные точки с подробными ответами}

\begin{enumerate}
    \item \textbf{Почему индукция не может быть обоснована логически?}\\
    Потому что из конечного множества единичных наблюдений нельзя логически вывести универсальное высказывание: всегда возможен новый факт, который ему противоречит.
    \item \textbf{Почему попытка обосновать индукцию опытом не помогает?}\\
    Потому что такое обоснование само будет индуктивным (``индукция работала раньше, значит, будет работать дальше''), что ведёт к кругу или бесконечному регрессу.
    \item \textbf{В чем заключается асимметрия подтверждения и опровержения?}\\
    Сколько угодно подтверждающих случаев не дают логического доказательства универсального закона, но один противоречащий случай (например, один чёрный лебедь против «все лебеди белые») логически несовместим с ним.
    \item \textbf{Что такое потенциальные фальсификаторы?}\\
    Это такие базисные эмпирические высказывания (описания наблюдаемых условий/событий), которые, если будут приняты, опровергнут теорию. Если у теории нет потенциальных фальсификаторов, то она не научна в попперовском смысле.
\end{enumerate}


\newpage
\section{Билет №18. Альтернативные классическому подходу теории истины. Проблема классификации теорий истины}

\subsection{Ответ}

Альтернативные корреспондентной теории истины возникают из недовольства тем, что в корреспондентной схеме остаются не вполне ясными два пункта: (i) как именно понимать «соответствие» и (ii) как мы вообще получаем доступ к реальности, с которой должно быть соответствие. В современной философии чаще всего обсуждают три альтернативные линии: \textbf{когерентную}, \textbf{прагматическую} и \textbf{дефляционную} (минималистскую).

\subsubsection*{1. Когерентная теория истины}

Когерентная концепция определяет истину как \textbf{системное свойство согласованности} убеждений: утверждение истинно, если оно «вписывается» в уже принятую систему без противоречий и поддерживает её целостность. Хинтикка отмечает, что такая трактовка типична для рационалистических построений, где акцент делается на внутреннюю связность теории.

Проблема, на которую (в духе критики Рассела) обычно указывают: возможны \textbf{разные} внутренне согласованные системы, которые несовместимы друг с другом. Тогда одной когерентности может не хватать, чтобы отличить истинное от просто «хорошо организованного».

\subsubsection*{2. Прагматическая теория истины}

Прагматическая линия (в частности, у Пирса) связывает истину не с «зеркальным соответствием», а с \textbf{успешностью практики и исследовательского процесса}. Важно понимать это не как грубую формулу «полезно значит истинно», а как идею о том, что истина мыслится через \textbf{долгосрочную перспективу исследования}: истинным считается то, к чему в пределе пришло бы рациональное сообщество исследователей при продолжении проверок и исправлений.

Отсюда часто говорят, что предикат «истинно» в прагматизме приближается к идее \textbf{оправданной утверждаемости}: мы называем «истинным» то, что выдерживает процедуры проверки и остаётся устойчивым в исследовательской практике.

Классическое возражение (которое разбирает Рассел): \textbf{полезность} и \textbf{истинность} не тождественны. Бывает полезная ошибка и вредная истина; поэтому прагматической теории требуется аккуратно объяснить, почему речь идёт не о сиюминутной выгоде, а об устойчивости и пределе исследования.

\subsubsection*{3. Дефляционные (минималистские) теории истины}

Дефляционные теории (Рамсей, Куайн) утверждают, что у истины нет «глубокой» метафизической природы: предикат «истинно» не обозначает особое свойство вещей или убеждений. Его роль в основном \textbf{языковая и логическая}.

Базовая интуиция здесь \textbf{дисквотационная} (``раскавычивание''): сказать
\begin{quote}
«Истинно, что идёт дождь»
\end{quote}
--- по содержанию то же самое, что сказать
\begin{quote}
«Идёт дождь».
\end{quote}
Именно поэтому Рамсей говорит об «избыточности» предиката истины: он как будто ничего не добавляет к содержанию конкретного утверждения.

При этом дефляционисты не обязательно отрицают полезность слова «истинно»: оно нужно как \textbf{инструмент} для обобщений и косвенного утверждения. Например, чтобы сказать «всё, что сказал свидетель, истинно», удобно иметь предикат истины, иначе пришлось бы перечислять все его фразы по одной. В этом смысле дефляционный подход пытается сохранить роль предиката истины, но без вывода, что истина --- отдельное «субстанциональное» свойство.

\subsubsection*{4. Проблема классификации теорий истины}

Проблема классификации, на которую указывает Анкин, состоит в том, что в учебной традиции часто смешивают разные вещи:
\begin{enumerate}
    \item \textbf{теорию истины} (что такое истина по своей природе),
    \item \textbf{критерии истины} (как мы распознаём истину на практике),
    \item \textbf{формальные требования адекватности} (например, совместимость с Т-схемой в тарсковском духе).
\end{enumerate}
Из-за этого одну и ту же позицию могут ошибочно записывать то в «теорию истины», то в «критерий истины». Анкин предлагает более фундаментальную ось: \textbf{инфляционизм} (истина как содержательное свойство, требующее философского объяснения) против \textbf{дефляционизма} (истина как логико-языковая функция без «глубокой» метафизики). На этой оси когерентизм и прагматизм обычно относятся к инфляционистским стратегиям, а минимализм/избыточность --- к дефляционным.

\subsection{Контрольные точки с подробными ответами}

\begin{enumerate}
    \item \textbf{В чём идея прагматической теории истины (Пирс) в минимально корректной формулировке?}\\
    Истина связывается с устойчивостью утверждения в практике исследования и с тем, к чему пришло бы сообщество исследователей в пределе проверки, а не просто с текущей «полезностью».
    \item \textbf{В чём суть дефляционной (минималистской) линии (Рамсей, Куайн)?}\\
    Предикат «истинно» не описывает особое свойство; в простых случаях он устраним через ``раскавычивание'', но остаётся полезным как логический инструмент для обобщений и косвенного утверждения.
    \item \textbf{В чём проблема классификации теорий истины (Анкин)?}\\
    Часто смешивают определение истины, критерии распознавания истины и формальные требования; поэтому корректнее делить подходы по оси инфляционизм/дефляционизм.
\end{enumerate}

\newpage
\section{Билет №19. Концепции знания. Классификации типов знания. Пропозициональное и предикативное знание}

\subsection{Ответ}

В аналитической философии под \textbf{концепциями знания} в данном билете понимаются разные теоретические схемы, которые по-разному описывают знание и делят его на виды: например, по способу доступа к объекту (Рассел: знакомство/описание) или по форме того, что мы знаем (Анкин; Райл: знание-что/знание-как/знание объектов). Поэтому здесь речь не об одном определении знания, а о нескольких классификационных подходах.

Как отмечает Гаспаров И. Г., наиболее значимым является различие между знанием-знакомством ($knowledge\ by\ acquaintance$) и знанием по описанию ($knowledge\ by\ description$), введённое Б. Расселом.

Знание-знакомство (предполагаемое как непосредственный когнитивный контакт субъекта с объектом) связано с тем, что объект дан субъекту прямо, без посредства описаний: Рассел относил сюда данные органов чувств, объекты памяти и акты интроспекции. Смысл этой категории в том, что она выступает фундаментом: это исходные «данности», на которые затем опирается построение более сложного знания. Однако такое знание ограничено личным опытом и не может быть передано другому иначе как через превращение в описание.

Знание по описанию --- это опосредованное знание о вещах, с которыми мы не знакомы непосредственно. Например, мы не знакомы с Юлием Цезарем, но знаем по описанию, что он был римским диктатором. Гаспаров подчеркивает, что знание по описанию опирается на знакомство с компонентами описания: мы понимаем описание благодаря знакомству с теми элементами (в том числе универсалиями), из которых оно составлено.

Современная классификация, представленная у Анкина Д. В. и в обсуждениях, связанных с Г. Райлом, выделяет три ключевых типа знания:
\begin{enumerate}
    \item \textbf{Знание-что} ($knowledge\ that$) --- \textbf{пропозициональное знание}. Оно имеет форму «$S$ знает, что $p$». Здесь $S$ --- субъект (носитель знания как состояния), а $p$ --- пропозиция, то есть содержание знания (то, что именно известно). Такое знание истинностно-значимо: пропозиция $p$ может быть истинной или ложной, и именно поэтому пропозициональное знание связывают с обоснованием.
    \item \textbf{Знание-как} ($knowledge\ how$) --- процедурное знание или знание-умение: навыки и способности (умение кататься на велосипеде, играть на скрипке). В аналитической традиции обсуждается, сводимо ли «знание-как» к «знанию-что»: интеллектуалисты трактуют умение как применение правил/знаний-фактов, а анти-интеллектуалисты (в рылевской линии) настаивают на автономности навыков.
    \item \textbf{Знание объектов (предикативное знание).} Оно фиксирует знание объекта через приписывание ему свойств/признаков и в современной логико-эпистемологической реконструкции часто стремится быть выраженным в пропозициональной форме: «знать объект» означает знать истинные пропозиции о его свойствах и отношениях (как отмечает Анкин). Поэтому предикативное знание нередко считают зависимым от пропозиционального.
\end{enumerate}

В результате пропозициональное знание обычно рассматривается как «знание в собственном смысле» для эпистемологии, поскольку оно напрямую связано с истинностью и обоснованием, тогда как знание-умение и знание объектов описывают другие важные, но не всегда сводимые формы знания.

\subsection{Контрольные точки с подробными ответами}

\begin{enumerate}
    \item \textbf{В чем разница между «знанием-что» и «знанием-как»?}\\
    «Знание-что» касается фактов и выражается в форме «$S$ знает, что $p$». «Знание-как» касается умений и проявляется в действии; спор состоит в том, сводится ли оно к набору пропозиций или является автономным.
    \item \textbf{В чем разница между знанием-знакомством и знанием по описанию у Рассела?}\\
    Знание-знакомство основано на непосредственной данности объекта, а знание по описанию относится к объектам без прямого знакомства и задаётся через описательные характеристики.
    \item \textbf{Почему предикативное знание часто считают зависимым от пропозиционального?}\\
    Потому что «знать объект» обычно разворачивают как «знать истинные утверждения о его свойствах и отношениях», то есть переводят в форму «$S$ знает, что $p$».
    \item \textbf{Что такое «универсалии» в концепции знания Рассела?}\\
    Это общие свойства/понятия (например, «белизна», «различие»), которые участвуют в описаниях и позволяют формулировать суждения.
\end{enumerate}


\newpage
\section{Билет №20. Истина и приближение к истине по К. Попперу. Правдоподобие и вероятность теорий}

\subsection{Ответ}

В поздних работах Поппер снимает своё прежнее подозрение к разговору об «абсолютной» истине и, опираясь на семантику Тарского, принимает корреспондентную установку: истина --- это соответствие фактам. Но для методологии науки одной этой идеи недостаточно, потому что на практике учёные часто работают с теориями, которые оказываются ложными или неточными. Поэтому Поппер вводит понятие \textbf{правдоподобия} (близости к истине): даже если теория не истинна строго, можно пытаться понять, стала ли новая теория \emph{ближе} к истине, чем старая.

\subsubsection*{1. Правдоподобие как баланс истинного и ложного содержания}

Поппер связывает правдоподобие с \textbf{содержанием} теории, разделяя его на истинное и ложное.
\begin{enumerate}
    \item \textbf{Истинное содержание} --- множество истинных следствий теории.
    \item \textbf{Ложное содержание} --- множество ложных следствий теории.
\end{enumerate}
Тогда теория $t_2$ более правдоподобна, чем $t_1$, если она улучшает баланс: у неё либо \emph{больше истинного содержания при том же (или меньшем) ложном}, либо \emph{меньше ложного при том же (или большем) истинном}. В этой логике наука может прогрессировать даже через смену опровергнутых теорий: новая теория может быть ложной в целом, но всё равно «лучше приближать» к истине.

\textbf{Уточнение к твоему вопросу про сравнение.} Если оказывается, что одна теория одновременно «даёт больше истины и больше лжи» (то есть растут и истинные, и ложные следствия), то по попперовскому критерию \emph{не всегда} получится однозначно сказать, что она более правдоподобна. В таком случае сравнение может быть \textbf{неопределённым} на уровне одной только формулы «истинное минус ложное»: теории могут быть просто несопоставимы по правилу доминирования (когда улучшение идёт по одной шкале, но ухудшение по другой). Тогда остаётся то, на чём Поппер настаивает методологически: сравнивать теории через \textbf{их проверяемость, рискованность предсказаний и результаты строгих испытаний} (то есть через то, какая теория лучше выдерживает критику и где именно обнаруживаются её ложные следствия).

\subsubsection*{2. Почему «вероятность теории» не является целью}

Важный пункт у Поппера --- противопоставление правдоподобия и вероятности. Он утверждает, что стремиться к «наиболее вероятным» теориям --- плохая методологическая цель, потому что высокая (логическая) вероятность обычно означает \textbf{бедное содержание}. Интуиция такая: чем больше теория утверждает о мире, тем больше она \textbf{запрещает} (то есть тем больше возможных наблюдений с ней несовместимы), а значит, тем \textbf{больше риск} оказаться ложной. Поэтому «смелые» универсальные теории имеют очень низкую логическую вероятность: они информативны именно потому, что легко могут быть опровергнуты.

Это можно уложить в попперовскую связку:
\begin{enumerate}
    \item высокая информативность $\Rightarrow$ высокая фальсифицируемость (много потенциальных фальсификаторов);
    \item высокая фальсифицируемость $\Rightarrow$ низкая (логическая) вероятность;
    \item следовательно, наука должна стремиться не к «вероятности», а к \textbf{содержательности и рискованным проверяемым предсказаниям}.
\end{enumerate}

\subsubsection*{3. Истина как регулятив и прогресс науки}

Для Поппера истина остаётся регулятивным идеалом: мы можем не знать, достигли ли истины, но можем критически улучшать теории, устраняя ошибки и увеличивая правдоподобие. Поэтому научный прогресс у него --- это движение к истине через рост содержательности и усиление проверок, а не через накопление подтверждений или поиск «наиболее вероятных» утверждений.

\subsection{Контрольные точки с подробными ответами}

\begin{enumerate}
    \item \textbf{А если $t_1$ даёт и больше истинных, и больше ложных следствий?}\\
    Тогда по одному лишь балансу «истинное/ложное содержание» сравнение может быть неразрешимым: критерий «$t_2$ лучше $t_1$» у Поппера требует улучшения без компенсирующего ухудшения (или наоборот). В таких смешанных случаях Поппер предлагает опираться на критические испытания: где именно теория проваливается, какие предсказания рискованнее и какая теория лучше выдерживает строгую проверку.
    \item \textbf{Почему Поппер против ориентации на вероятность теорий?}\\
    Потому что «вероятнее» обычно значит «менее содержательно»: теория с высоким содержанием больше запрещает и потому более рискованна, а значит, имеет низкую логическую вероятность. Наука, по Попперу, ценит не вероятные, а смелые и проверяемые теории, которые можно подвергнуть фальсификации.
\end{enumerate}

\newpage
\section{Билет №21. Проблема обоснования, современные концепции и проблемы. Фундаментализм и антифундаментализм. Трилемма Мюнхгаузена}

\subsection{Ответ}

Проблема обоснования является центральной для эпистемологии, так как именно обоснование ($justification$) отделяет знание от случайно истинного мнения. Как отмечает Гаспаров И. Г., в современной философии под обоснованием понимают наличие у субъекта достаточных причин считать свое убеждение истинным.

Главное методологическое препятствие здесь — Трилемма Мюнхгаузена, сформулированная Гансом Альбертом. Пытаясь обосновать утверждение $A$, мы ссылаемся на основание $B$. Но $B$ само требует обоснования. Это ведет к трем неприемлемым вариантам:
1. Бесконечный регресс: процесс обоснования никогда не заканчивается, и мы не достигаем фундамента.
2. Логический круг: утверждение обосновывается через то, что само зависит от этого утверждения.
3. Догматизм (прерыв обоснования): мы произвольно выбираем точку, которая не требует обоснования (догму).

По отношению к этой трилемме выделяются две основные позиции: фундаментализм и антифундаментализм (когерентизм и инфинитизм).

Фундаментализм настаивает на существовании «эпистемического базиса». Согласно Гаспарову, фундаменталисты выбирают третий путь трилеммы, но отрицают, что это догматизм. Они утверждают, что существуют самоочевидные, базовые убеждения, которые не нуждаются в обосновании другими убеждениями. Классический пример — когито Декарта или чувственные данные (эмпиризм). Гаспаров указывает на проблему фундаментализма: как именно базовые убеждения передают свою обоснованность небазисным? Логический вывод сохраняет истину, но не всегда «понятость» обоснования.

Антифундаментализм представлен прежде всего когерентизмом. По Гаспарову, когерентисты выбирают второй путь трилеммы — круг, но называют его «сетью». Обоснованность убеждения зависит не от фундамента, а от его согласованности ($coherence$) со всей системой убеждений. Если убеждение гармонично вписывается в общую картину мира, оно обосновано. Главная проблема когерентизма («проблема изоляции») — система может быть идеально согласованной внутри (как хороший роман), но полностью оторванной от реальности.

Отдельно Гаспаров выделяет релайабилизм как современную экстерналистскую концепцию. Здесь обоснование заменяется надежностью когнитивного процесса. Если убеждение сформировано механизмом, который дает высокий процент истинных результатов (зрение, память), оно обосновано, даже если субъект не может привести логических доводов «изнутри».

Принцип «достаточного основания» Лейбница, согласно анализу Альберта, в классической науке служил постулатом рациональности, требующим искать причину для каждого факта. Однако Альберт доказывает, что этот принцип недостижим в абсолютном смысле и ведет к признанию погрешимости любого разума.

\subsection{Обзор источников}

Основной философский каркас темы задает \textbf{Ганс Альберт в «Трактате о критическом разуме»} (гл. 1). Он вводит понятие «Трилемма Мюнхгаузена» и критикует «принцип достаточного основания» как скрытый догматизм. \textbf{Гаспаров И. Г. «Knowledge»} (раздел 9.2) систематизирует ответы на трилемму Альберта, классифицируя их на фундаментализм, когерентизм и инфинитизм. \textbf{Лебедев и Черняк} дополняют этот анализ спором интернализма и экстернализма, объясняя, почему современные концепции (релайабилизм) пытаются уйти от «чистой логики» к психологии познания.

\subsection{Контрольные точки с подробными ответами}

\textbf{1. В чем заключается суть Трилеммы Мюнхгаузена?}\\
Это невозможность построить абсолютно надежное обоснование без того, чтобы либо уйти в бесконечность, либо ходить по кругу, либо остановиться на бездоказательном утверждении.

\textbf{2. Чем базовые убеждения отличаются от небазисных в фундаментализме?}\\
Базовые убеждения самоочевидны или самообоснованы (например, боль), небазисные получают свою силу только через логическую связь с базовыми.

\textbf{3. Какова главная претензия Альберта к принципу достаточного основания?}\\
Альберт считает, что это требование ведет к интеллектуальному тупику, так как поиск «окончательного» основания всегда иллюзорен и заканчивается догматизмом.

\textbf{4. Что такое «проблема изоляции» в когерентизме?}\\
Это риск того, что система убеждений будет логически стройной и согласованной, но при этом ложной, так как когерентность проверяет связь между мыслями, а не связь мыслей с миром.



\newpage
\section{Билет №22. Проблема обоснования у К. Поппера – Введение к «Предположениям и опровержениям» («Об источниках знания и невежества»)}

\subsection{Ответ}

Во вводном тексте «Об источниках знания и невежества» Поппер пересматривает саму постановку проблемы обоснования. Традиционная эпистемология, по Попперу, задавала вопрос в форме: \emph{каков главный источник истинного знания} — чувства (эмпиризм) или разум (рационализм). Поппер называет такой вопрос \textbf{авторитарным}, потому что он подразумевает поиск привилегированного и в идеале непогрешимого «суда», который должен \emph{легитимизировать} знание через его происхождение (как будто правильное происхождение автоматически гарантирует истинность результата).

Поппер отвергает эту «источник-ориентированную» модель. Его тезис: \textbf{не существует идеальных источников знания}. И чувственные данные, и рассуждение могут вести к ошибкам. Если же пытаться построить знание на «позитивном обосновании», то мы сталкиваемся с классической дилеммой:
\begin{enumerate}
    \item либо мы останавливаемся на некотором основании как на безусловно надёжном (догматизм),
    \item либо требуем основания для основания и уходим в бесконечный регресс.
\end{enumerate}
Поэтому Поппер предлагает заменить вопрос «\emph{откуда} берётся знание?» на вопрос «\emph{как} обнаруживать и устранять ошибки?». Это переход от генетического/авторитарного подхода к \textbf{критическому}.

Центральный мотив текста --- критика \textbf{учения об очевидности истины} (эпистемологического оптимизма). Поппер обсуждает классические варианты этой установки: у Декарта надежда на метод, ведущий к очевидности (подкреплённой гарантией истинности), у Бэкона надежда на очищенный опыт как «правдивый» источник. Поппер подчёркивает, что такая вера опасна: если истина якобы сама себя проявляет, то несогласие легко объяснить не сложностью поиска истины, а «слепотой» людей, которую кто-то должен вызвать. В результате появляется склонность к конспирологическим объяснениям и к нетерпимости: вместо признания собственной ошибочности удобнее объявить оппонента неразумным или злонамеренным.

Альтернатива Поппера --- \textbf{критический рационализм}. Его главный принцип в контексте обоснования: мы не должны пытаться \emph{доказывать} истинность теорий, мы должны пытаться \emph{критиковать} их и выявлять ошибки. Теории остаются \textbf{догадками} (предположениями), а рациональность обеспечивается процедурой их проверки и готовностью отказаться от теории при серьёзных возражениях. В этом смысле «обоснование» заменяется \textbf{подкреплением} (corroboration): теория считается предпочтительной не потому, что получила окончательное оправдание, а потому, что выдержала строгие попытки опровержения лучше конкурентов.

При этом Поппер подчёркивает роль \textbf{традиции}: практически всё наше знание приходит к нам через уже существующие теории, языковые практики и научные школы; мы не начинаем «с нуля». Но традиция не должна превращаться в авторитет: её ценность именно в том, что она даёт материал для критики и улучшения. Поэтому прогресс знания у Поппера --- это не построение на несомненном фундаменте, а бесконечный процесс \textbf{исправления ошибок}.

\subsection{Контрольные точки с подробными ответами}

\begin{enumerate}
    \item \textbf{Почему Поппер считает вопрос об источниках знания авторитарным?}\\
    Потому что он ищет привилегированный «гарантирующий» авторитет (чувства, разум и т.п.), который должен оправдать знание по происхождению, вместо того чтобы оценивать теории по критике и испытаниям.
    \item \textbf{Почему, по Попперу, «позитивное обоснование» ведёт к догматизму или регрессу?}\\
    Любая попытка окончательно оправдать знание требует либо принять какое-то основание как непогрешимое (догматизм), либо бесконечно требовать основания для основания (регресс).
    \item \textbf{К каким последствиям ведёт вера в очевидность истины?}\\
    Она делает несогласие «подозрительным» и провоцирует объяснять ошибки не трудностью поиска истины, а чьей-то слепотой/злонамеренностью, что подпитывает нетерпимость и конспирологические модели.
    \item \textbf{Как критический рационализм заменяет традиционную проблему обоснования?}\\
    Он заменяет вопрос «почему мы вправе верить?» на вопрос «как найти ошибку?»: теории выдвигаются как догадки и оцениваются через критику, проверки и попытки опровержения; выжившие теории считаются подкреплёнными, но не доказанными.
    \item \textbf{Какова роль традиции у Поппера?}\\
    Традиция --- главный источник содержания наших знаний, но рациональность требует не подчинения традиции, а её постоянной критики и улучшения.
\end{enumerate}

\end{document}
